% ============================================================================
%  PART II: How SiliconScribe works
% ============================================================================

\chapter{Architecture and Code Map}
\label{ch:arch}

\section{Three tiers}

\begin{figure}[H]
\centering
\begin{tikzpicture}[node distance=10mm and 8mm]
  \node[hbox,minimum width=13.4cm] (fe) {FRONTEND \quad React 18 + Vite + TypeScript (runs in the browser)};
  \node[box,below=3mm of fe.south west,anchor=north west,minimum width=13.4cm,align=left,font=\scriptsize] (fed)
    {Sidebar: prompt, run timeline, examples, Model menu, ``Testbench written by'', specification review, target clock \\
     ``Ask about this run'' (the debugging agent) under the run timeline \\
     Design / Results / Test cases / Waveform / Schematic tabs \textbullet\ Projects menu \textbullet\ Export UVM, Re-run, Stop \textbullet\ theme};
  \node[hbox,below=12mm of fed.south west,anchor=north west,minimum width=9.1cm] (be) {BACKEND \quad FastAPI + Pydantic (Python)};
  \node[box,below=3mm of be.south west,anchor=north west,minimum width=9.1cm,text width=8.8cm,align=left,font=\scriptsize] (bed)
    {\file{main.py}: routes, CORS, rate limit, SSE streaming\\
     \file{orchestrator.py}: pipeline, repair loop, second opinion\\
     \file{llm_service.py} + \file{model_catalog.py}: providers, model menu, prompts\\
     \file{simulator.py}: Icarus Verilog and sandbox\\
     \file{mutation.py}: testbench strength \textbullet\ \file{timing.py}: timing and area\\
     \file{equivalence.py}: proof that the netlist equals the RTL\\
     \file{agent.py} + \file{wave_query.py}: the debugging agent\\
     linter, checks, VCD parser and toggle coverage, schematic, synthesis, UVM export, SQLite};
  \node[hbox,right=4mm of be,minimum width=3.9cm] (llm) {LLM PROVIDERS};
  \node[box,below=3mm of llm.south west,anchor=north west,minimum width=3.9cm,text width=3.6cm,align=left,font=\scriptsize] (llmd)
    {OpenCode app (local)\\ Cloudflare Workers AI\\ NVIDIA NIM\\ any OpenAI-compatible API\\ keys only in \file{backend/.env}};
  \node[hbox,below=12mm of bed.south west,anchor=north west,minimum width=9.1cm] (eda) {EDA TOOLCHAIN \quad local subprocesses};
  \node[box,below=3mm of eda.south west,anchor=north west,minimum width=9.1cm,align=left,font=\scriptsize] (edad)
    {Icarus Verilog 12 (\code{iverilog}, \code{vvp}) \textbullet\ Yosys 0.71 \textbullet\ OpenSTA 3.2.0 \textbullet\ sky130 cell library\\ Verilator if installed \textbullet\ one workspace folder per design, time limits, scrubbed environment};
  \draw[arr,<->] (fed.south) ++(-2.5,0) -- node[right,font=\scriptsize]{HTTP + Server-Sent Events} ++(0,-1.2);
  \draw[arr,<->] (bed.east) -- (llmd.west);
  \draw[arr] (bed.south) ++(-1,0) -- node[right,font=\scriptsize]{subprocess} ++(0,-1.2);
\end{tikzpicture}
\caption{The three tiers. The backend is the only part that holds secrets or runs programs.}
\end{figure}

\begin{keyidea}
The browser never talks to the AI provider and never runs tools. All secrets (API keys) and all execution stay in the backend, which is why the frontend can be a static website. The web page never asks for a key: keys are read only from \file{backend/.env} (or the environment of the terminal that starts the backend).
\end{keyidea}

\section{Folder map}

\begin{center}
\small
\begin{longtable}{@{}p{5.3cm}p{9.8cm}@{}}
\toprule
\textbf{Path} & \textbf{What lives there} \\ \midrule \endhead
\file{backend/main.py} & The web server: every HTTP endpoint, the rate limiter, CORS, SSE streaming, saving finished runs. \\
\file{backend/orchestrator.py} & The ``brain'': runs the stages in order, validates outputs, runs the self-correction loop and its guards, re-runs the original testbench, asks for a second opinion, starts the strength and timing checks, handles fallback. \\
\file{backend/llm_service.py} & Talks to AI providers: the system prompts, provider selection (OpenCode app, Cloudflare, NVIDIA NIM, OpenAI-compatible, offline), per-stage model routing, response parsing. \\
\file{backend/model_catalog.py} & Builds the Model menu: reads the OpenCode app's model list and the live Cloudflare and NVIDIA NIM catalogues, caches them, and checks every model id before it is used. \\
\file{backend/simulator.py} & Runs Icarus Verilog safely: sandbox policy, file writing, compile, run, verdict. \\
\file{backend/mutation.py} & Testbench strength: injects small bugs into a passing design and counts how many the testbench catches. \\
\file{backend/timing.py} & Pre-layout timing and area: Yosys maps the design to sky130 cells, OpenSTA times the netlist and gives the critical path pin by pin. \\
\file{backend/equivalence.py} & Formal equivalence: proves with Yosys' built-in SAT solver that the sky130 netlist computes the same function as the RTL. \\
\file{backend/agent.py} & The debugging agent behind ``Ask about this run'': the JSON tool loop, its eight tools, fix verification and the rule-based fallback (Chapter~\ref{ch:agent}). \\
\file{backend/wave_query.py} & \code{WaveDB}: a queryable copy of one run's VCD, with values at and just before any time, edges, changes and value search. \\
\file{backend/pdk/} & The sky130 cell library file (\file{sky130hd_tt.lib}) that timing needs, with a README saying where to fetch it. \\
\file{backend/models.py} & The data contracts (Pydantic models) for every request, response and record. \\
\file{backend/offline_designs.py} & Built-in designs (counter, 4- and 8-bit ALU, full adder, 4:1 mux) with their testbenches, used offline and for the mutation sweep. \\
\file{backend/linter.py} & Non-blocking lint: Verilator, then Icarus, then pattern heuristics. \\
\file{backend/cov_metrics.py} & Builds the coverage summary from the simulation result. \\
\file{backend/checks.py} & Turns the testbench's \code{CHK|...} lines into per-check rows and exact per-type counters for the Test cases tab (display only; never decides the verdict). \\
\file{backend/vcd_parser.py} & Reads the VCD waveform into JSON for the browser, and measures toggle coverage from the same file. \\
\file{backend/schematic.py}, \file{synthesis.py} & Port-level block diagram; generic (unmapped) Yosys synthesis with honest labelling, which also produces the two reduced netlists for the RTL-cells and gate-level schematic views. \\
\file{backend/uvm_templates.py} & Generates the 16-file UVM export bundle. \\
\file{backend/db.py}, \file{cleanup.py} & SQLite project storage; deletion of old workspace folders. \\
\file{backend/openroad.py} & A stub for a future physical-design flow (reports ``not available''). \\
\file{backend/tools/mutation_sweep.py} & The 27-mutant testbench-quality sweep over the built-in designs. \\
\file{backend/tools/verify_mapping.py} & Checks the timing stage itself: simulates the mapped gate netlist with the same testbench, recounts area and cells, and checks that the minimum period does not depend on the target. \\
\file{backend/tools/preflight.py} & Demo-day readiness check: tools, library, models, offline demo, saved results (with \code{--live}, one tiny model call). \\
\file{backend/tools/evaluate_live.py} & Runs the fixed 15-prompt live evaluation and writes \file{evaluation/results.jsonl} and \file{summary.md}. \\
\file{backend/tools/build_opensta.sh} & Builds OpenSTA from source into \file{~/.local/opensta} (no Homebrew). \\
\file{backend/evaluation/} & Recorded live results: \file{results.jsonl}, \file{summary.md}, the final designs and testbenches in \file{artifacts/}. \\
\file{backend/evaluation/independent/} & Hand-written reference testbenches (\file{tb_c1.v} \ldots\ \file{tb_s6.v}), \file{run.py} that checks each evaluated design against them, and \file{strength.py} that checks the references themselves with every mutant. \\
\file{backend/tests/} & 632 automated backend tests in 36 test files, plus \file{conftest.py} (shared test set-up) and two fake servers (\file{cf_fake.py}, \file{nim_fake.py}) so Cloudflare and NIM code is tested without the internet. \\
\file{backend/workspace/} & Runtime files: one folder per design (\file{design.sv}, \file{testbench.sv}, \file{transcript.log}, \file{design.vcd}) and \file{projects.db}. \\
\file{frontend/src/App.tsx} & The page: holds all state and connects the components. \\
\file{frontend/src/api.ts} & Every call to the backend, including the SSE stream reader. \\
\file{frontend/src/types.ts} & TypeScript copies of the backend data shapes. \\
\file{frontend/src/theme.ts} & Light, dark or system theme; an explicit choice is remembered in the browser. \\
\file{frontend/src/components/} & PromptPanel, ModelPicker, TestbenchModelSelect, SpecReview, SpecSummary, AgentChat, DebugChat, CodeEditor, ResultsPanel, TestCasesView, WaveformViewer, SchematicView (plus their tests). \\
\file{frontend/src/ErrorBoundary.tsx} & Catches a drawing error in the results panel so the rest of the page keeps working. \\
\file{frontend/src/netlistRender.ts} & Lays out a Yosys netlist as an SVG drawing with netlistsvg; loaded only when a schematic view needs it. \\
\file{docs/} & Report, presentation, poster, screenshots, live evidence, this guide. \\
\file{run.sh}, \file{Dockerfile}, \file{docker-compose.yml}, \file{render.yaml}, \file{vercel.json} & Ways to run and deploy (Chapter~\ref{ch:running}). \\
\bottomrule
\end{longtable}
\end{center}

% ---------------------------------------------------------------------------
\chapter{The Journey of One Request}
\label{ch:journey}

This chapter follows a single prompt from the moment you press Generate RTL to the final PASS. Each step names the function that does the work, so you can open the code and follow along.

\begin{figure}[H]
\centering
\newcommand{\msg}[5]{\draw[arr#5] (#1,#3) -- node[above,font=\scriptsize]{#4} (#2,#3);}
\resizebox{\textwidth}{!}{%
\begin{tikzpicture}[font=\scriptsize]
  \foreach \x/\n/\l in {0/B/Browser,3.1/A/FastAPI route,6.2/O/Orchestrator,9.3/L/LLM provider,12.4/S/EDA tools}{
    \node[hbox,minimum width=2.6cm] (\n) at (\x,0) {\l};
    \draw[midgrey,dashed] (\x,-0.45) -- (\x,-10.4);
  }
  \msg{0}{3.1}{-1.0}{POST /api/design/stream}{}
  \msg{3.1}{6.2}{-1.6}{pipeline\_events()}{}
  \msg{6.2}{9.3}{-2.2}{parse\_intent}{}
  \msg{9.3}{6.2}{-2.7}{RTLDesignSpec}{,dashed}
  \msg{6.2}{0}{-3.1}{SSE: start, intent}{,dashed}
  \msg{6.2}{9.3}{-3.7}{generate\_rtl, generate\_testbench, explain}{}
  \msg{6.2}{0}{-4.3}{SSE: rtl, lint, testbench, explanation}{,dashed}
  \msg{6.2}{12.4}{-4.9}{simulate(): iverilog + vvp}{}
  \msg{12.4}{6.2}{-5.4}{transcript + VCD}{,dashed}
  \msg{6.2}{0}{-5.8}{SSE: simulate (PASS / FAIL / ERROR)}{,dashed}
  \node[draw=badred,dashed,rounded corners,minimum width=10.2cm,minimum height=2.2cm,anchor=north west] (loopbox) at (5.6,-6.15) {};
  \node[anchor=north east,font=\scriptsize,text=badred] at ($(loopbox.north east)+(-0.1,-0.08)$) {loop while FAIL/ERROR/TIMEOUT and iterations remain};
  \msg{6.2}{9.3}{-7.0}{fix\_design (or compiler repair first)}{}
  \msg{6.2}{0}{-7.5}{SSE: fixing, fix}{,dashed}
  \msg{6.2}{12.4}{-8.0}{re-simulate}{}
  \msg{6.2}{12.4}{-9.05}{after a PASS: mutants, Yosys + OpenSTA}{}
  \msg{6.2}{0}{-9.55}{SSE: synthesis, strength, timing, done}{,dashed}
  \node[font=\scriptsize,align=left,text=inkgrey] at (3.1,-10.1) {route saves the project to SQLite};
\end{tikzpicture}}
\caption{Sequence of one streamed run. Solid arrows are calls; dashed arrows are replies and streamed events.}
\label{fig:sequence}
\end{figure}

\section{Step by step}

\begin{enumerate}[leftmargin=*,label=\textbf{Step \arabic*.}]
  \item \textbf{You type and press Generate RTL} (or Ctrl/Cmd+Enter). \code{PromptPanel.submit} rejects an empty prompt or one longer than 2000 characters, clamps the target clock to 1--10000 MHz (the slider itself goes from 10 to 500) and calls \code{App.handleGenerate}.
  \item \textbf{Optional: review the specification.} If ``Review the specification before generating'' is ticked, \code{handleGenerate} first calls \code{api.parseSpec}, which sends \code{POST /api/spec/parse}. The server only runs \code{parse_intent} and returns the specification. The page shows it in the SpecReview panel, you edit the behaviour and constraints, and ``Generate with this specification'' starts the run with that exact specification (Section~\ref{sec:specreview}).
  \item \textbf{The browser opens a stream.} \code{startRun} clears the previous results and calls \code{api.streamDesign}, which sends \code{POST /api/design/stream} with the prompt, target clock, the chosen model, the optional testbench model, the optional confirmed specification, the thinking level, \code{self_correct=true}, \code{max_iterations=5} and \code{timeout_seconds=30}.
  \item \textbf{The server checks the request.} FastAPI validates the JSON against \code{RunRequest} (HTTP 422 if invalid). This includes the model ids (Section~\ref{sec:modelid}) and the bounds on a user specification. \code{design_stream} applies the rate limit, creates a fresh design id with \code{_new_design_id} (a full random UUID) and starts \code{pipeline_events}.
  \item \textbf{Specification.} The orchestrator calls \code{parse_intent}, unless you already confirmed a specification. The model returns JSON describing the module name, parameters, ports (name, direction, width), behaviour and constraints, which is validated as an \code{RTLDesignSpec}. Event \code{intent} is streamed.
  \item \textbf{RTL.} \code{generate_rtl} is called with the specification. The orchestrator checks the reply contains a complete \code{module <name> ... endmodule}; if not, it retries once, then fails clearly. Event \code{rtl}. The linter may add a \code{lint} event.
  \item \textbf{Testbench.} \code{generate_testbench} is called with the specification (by the separate testbench model if you chose one). The result must be a complete module named \code{testbench}, must really instantiate the design and must contain at least one check; otherwise retry once, then refuse. Event \code{testbench}.
  \item \textbf{Explanation.} \code{explain_design} writes 2--3 sentences. If this cosmetic step fails, a placeholder is used and the run continues. Event \code{explanation}.
  \item \textbf{First simulation.} \code{simulator.simulate} scans the code against the sandbox policy, writes the two files, compiles, runs and parses the transcript. The result is recorded as iteration 0. Event \code{simulate}.
  \item \textbf{Self-correction loop} (only if FAIL, ERROR or TIMEOUT): compiler repair or LLM diagnosis, guards, re-simulation, up to \code{max_iterations} times (Chapter~\ref{ch:pipeline}). Events \code{fixing}, \code{fix}, \code{simulate}.
  \item \textbf{Post-processing.} Coverage summary, VCD parsed into a waveform, port-level schematic, generic synthesis (event \code{synthesis}, which also carries the two netlists the Schematic tab draws). If the run passed only after the testbench was changed, the original testbench is run again on the final design, and a fresh second-opinion testbench is requested (event \code{second_opinion}). Then, for a PASS only, testbench strength (event \code{strength}) and the timing and area estimate with its critical path and equivalence proof (event \code{timing}). Chapters~\ref{ch:trust} and~\ref{ch:timing} explain these.
  \item \textbf{Done.} A full \code{RunResponse} is built and sent in the \code{done} event. It records which model wrote the design and, if different, which wrote the testbench. \code{design_stream} saves it to SQLite before forwarding it.
  \item \textbf{The page updates.} For each event, \code{startRun} adds a message to the run timeline in the sidebar and fills the editors and result panels; on \code{done} it loads the final response, switches to the Results tab and refreshes the Projects list.
\end{enumerate}

\section{The event protocol}

Every streamed message is one JSON object with a \code{stage} field. The full list:

\begin{center}
\small
\begin{tabularx}{\textwidth}{@{}>{\ttfamily}lX@{}}
\toprule
\textrm{\textbf{stage}} & \textbf{Meaning and extra fields} \\
\midrule
start & Run started; names the provider and model. \\
intent & Specification ready: \code{rtl_spec}, and \code{fallback_notice} if degraded. \\
rtl & RTL ready: \code{rtl_code}. \\
lint & Lint found something: \code{target} (rtl, tb or fix) and \code{lint} (errors, warnings). Informational only. \\
testbench & Testbench ready: \code{testbench_code}. \\
explanation & \code{explanation} text. \\
simulate & One simulation finished: \code{iteration}, \code{status}, \code{result}. \\
fixing & About to ask the model to diagnose (iteration number). \\
fix & A fix was applied, or the loop stopped: \code{fix_summary}, \code{fix_type} (compiler\_repair, llm\_diagnosis or offline\_scripted), \code{diagnosis}, new code. \\
synthesis & Generic synthesis result or ``not available''. \\
second\_opinion & Sent twice when a repair rewrote the testbench: first ``asking for a fresh testbench'', then \code{second_opinion} (status PASS, FAIL or unavailable, with counts and the model). \\
strength & After a PASS: \code{strength} (bugs injected, bugs caught, survivors). \\
timing & After a PASS, when the tools exist: \code{timing} (slack, area, cells, or the reason it was refused). \\
done & Final verdict: \code{status}, \code{response} (the complete RunResponse, including \code{pass_disputed} when the evidence contradicts a PASS); its message says ``PASS, but disputed: \ldots'' in that case. \\
error & Something failed: a short, safe \code{message}. \\
\bottomrule
\end{tabularx}
\end{center}

% ---------------------------------------------------------------------------
\chapter{The Pipeline and Self-Correction Loop}
\label{ch:pipeline}

All of this lives in \file{backend/orchestrator.py}, mainly in the generator function \code{pipeline_events}. A \textbf{generator} is a Python function that produces values one at a time with \code{yield}; each yielded value here is one event, which is exactly what streaming needs. The one-shot route uses \code{run_pipeline}, which simply consumes the same events and returns the final response.

\section{Validation between stages}

The orchestrator does not pass model output blindly to the next stage:

\begin{description}
  \item[\code{_has_complete_module(code, name)}] True only if the text contains \code{endmodule} and a \code{module <name>} header, both outside comments and strings (a comment such as \code{// module counter} next to \code{module up_counter} used to pass). Catches truncated or chatty answers before they reach the compiler, where they would cause a cryptic syntax error.
  \item[\code{_tb_check_count(tb)}] Counts the places where a testbench can actually fail: increments of a fail or error counter, \code{$error} / \code{$fatal}, \code{assert(} and calls of the standard \code{chk(...)} task. Comments and strings are blanked out first (\code{_strip_comments_and_strings}), so a comment such as \code{// $error} does not count, and a bare \code{fail_count = 0} is not a check. A testbench with zero checks cannot fail, so it is refused: ``Refusing to report PASS from a testbench that cannot fail.''
  \item[\code{_tb_problems(tb, name)}] The full acceptance test for a testbench: a complete module named \code{testbench}, at least one check, and a real instance of the design (\code{_instantiates_dut}), not just a mention of its name. The same rules apply to a testbench proposed by a repair.
  \item[Retry once.] RTL and testbench generation each get one automatic retry when the answer is invalid; a second failure stops the run with a clear message.
\end{description}

\section{The self-correction loop}

\begin{figure}[H]
\centering
\begin{tikzpicture}[node distance=7mm and 10mm]
  \node[box] (sim) {Simulate};
  \node[diamond,draw=navy,align=center,aspect=2.2,inner sep=1pt,font=\small,below=of sim] (d1) {PASS?};
  \node[okbox,right=of d1] (done) {Done};
  \node[diamond,draw=navy,align=center,aspect=2.2,inner sep=1pt,font=\small,below=of d1] (d2) {ERROR with a\\ known compiler\\ message?};
  \node[box,right=of d2] (cr) {Compiler repair\\ \scriptsize wire$\to$reg, disable$\to$return};
  \node[box,below=of d2] (llm) {LLM diagnosis + patch\\ \scriptsize symptom, root cause, fix, confidence};
  \node[box,below=of llm] (ac) {Anti-cheat: drop TB change\\ if it removed checks};
  \node[diamond,draw=navy,align=center,aspect=2.2,inner sep=1pt,font=\small,below=of ac] (d3) {Changed and\\ never seen?};
  \node[badbox,right=of d3] (stop) {Stop: no-op\\ or oscillation};
  \draw[arr] (sim)--(d1); \draw[arr] (d1)--node[above,font=\scriptsize]{yes}(done);
  \draw[arr] (d1)--node[left,font=\scriptsize]{no}(d2);
  \draw[arr] (d2)--node[above,font=\scriptsize]{yes}(cr);
  \draw[arr] (d2)--node[left,font=\scriptsize]{no}(llm);
  \draw[arr] (llm)--(ac); \draw[arr] (ac)--(d3);
  \draw[arr] (d3)--node[above,font=\scriptsize]{no}(stop);
  \draw[arr] (cr.north) |- (sim.east);
  \draw[arr] (d3.west) -- ++(-1.6,0) |- node[pos=0.25,left,font=\scriptsize]{yes: re-simulate} (sim.west);
\end{tikzpicture}
\caption{One pass of the loop. It also stops when \code{max_iterations} is reached (5 from the UI).}
\end{figure}

The loop condition is:
\begin{lstlisting}[style=plain]
while self_correct and result.status in ("FAIL", "ERROR", "TIMEOUT") and iteration < max_iterations:
\end{lstlisting}
Each pass is recorded as an \code{IterationRecord} (iteration number, status, fix summary, fix type, diagnosis, pass/fail counts, last 1200 characters of the log) so the UI can show a timeline.

\subsection{Deterministic compiler repairs}
\label{sec:compilerrepair}

Some failures are mechanical, and the compiler names the exact problem. Asking an LLM to guess would be slower and less reliable, so \code{_compiler_repairs} fixes two of them directly:

\begin{description}
  \item[\code{_repair_compile_errors} (wire to reg).] Icarus reports \code{design.sv:LINE: error: X is not a valid l-value} when a \code{wire} is assigned inside an \code{initial} or \code{always} block. The function finds the file the compiler blamed and changes only that declaration, \code{wire [3:0] X;} to \code{reg [3:0] X;}.
  \item[\code{_repair_disable_functions} (disable to return).] Icarus rejects \code{disable fname;} used as an early exit inside a function. It is replaced with \code{return fname;}, which Icarus accepts and which keeps the first-match behaviour.
\end{description}

Each repair produces a structured diagnosis with confidence ``high'' and the compiler line as evidence, and is labelled \code{fix_type = "compiler_repair"} so it is never confused with AI reasoning.

\subsection{LLM diagnosis}

Otherwise, \code{fix_design} sends the debugger model the specification (``source of truth''), the RTL, the testbench and the last 4000 characters of the simulator log. The prompt requires the answer in tagged sections:

\begin{lstlisting}[style=plain]
<diagnosis>
symptom: ...        root_cause: ...
fix_location: ...   confidence: high|medium|low
</diagnosis>
<rtl> ...full corrected RTL... </rtl>
<testbench> ...full testbench... </testbench>
<summary> one sentence </summary>
\end{lstlisting}

The prompt also tells the model to check the testbench against the specification first: if the testbench's expected value is wrong, the bug is in the testbench, not the RTL. This rule came from a real live failure (a 32-bit expected value compared with a 3-bit output).

Two things can go wrong with the answer itself, and both are handled honestly. If the reply cannot be read (a refusal, a reply cut off inside \code{<rtl>}, code without the tags), the fixer asks once more with the format spelled out; if the second reply is also unreadable the loop stops with ``The repair reply could not be read'', not with the misleading ``no change''. And if the model stops answering in the middle of the loop (a quota or a timeout), the run is \emph{kept}: the last simulated result, the code and the waveform are returned with the note ``The repair model became unavailable \ldots{} this is the last simulated result''. Before the last review such a failure threw the whole run away.

\subsection{Anti-cheat guard}

The testbench is the judge, so the repair loop must not be able to weaken it. The rule depends on what the last run showed:

\begin{description}
  \item[The testbench did not compile.] Its oracle never ran, so the fixer may rewrite it freely.
  \item[It ran but printed no readable summary] (for example \code{passed = 12} instead of \code{Passed: 12}) and no failure lines. The oracle is broken, so it may be repaired, but the repaired testbench may not have fewer checks.
  \item[Anything else] (a FAIL, a runtime stop such as \code{$fatal}, a hang). The testbench did its job. A proposed testbench with \emph{fewer} checks is thrown away with the note ``Testbench change rejected: it removed self-checks''.
\end{description}

The middle case used to be treated like the first: any ERROR or TIMEOUT unlocked a free rewrite. The last review showed how that could be abused: a testbench that stopped at its first mismatch with \code{$fatal} gave an ERROR, the design model replaced it with a one-check testbench, kept the bug, and got a PASS. Failure lines without a summary are now a FAIL (Section~\ref{sec:verdict}), so that path is closed. A proposed RTL that is not a complete module with the right name, or a proposed testbench that fails \code{_tb_problems}, is also rejected.

When a separate model wrote the testbench (Section~\ref{sec:independenttb}), the rule is stricter: the design model may not change that testbench at all unless its oracle is broken (one of the first two cases). If it does repair it, the result is no longer labelled independent, and a warning says so (``The separate model's testbench did not compile and was repaired by the design model, so the final testbench is no longer independent of the design'').

Even an accepted testbench change leaves a trace. At the end of the run, if the result is PASS and the testbench is no longer the one first generated, the \emph{original} testbench is run again on the final design. If it still fails, the result carries the warning ``The testbench was modified during repair, and the ORIGINAL testbench still reports \ldots{} Review the testbench change before trusting this PASS.'' The verdict itself then changes too: the response gets a \code{pass_disputed} reason and the screen shows \textbf{PASS, disputed} in warning colour with ``Do not trust this PASS yet'' under it. The same happens when the fresh second-opinion testbench fails the final design. Before the last review the headline stayed a plain green PASS and the contradiction was only in the warning text, which is exactly where an examiner would not look.''

\subsection{No-op and oscillation guards}

\code{_state_hash} computes a SHA-256 fingerprint of the RTL and testbench together, after removing comments and all whitespace, so cosmetic edits do not count as changes. If the fix produces the same fingerprint as now, the loop stops (``Model returned no change''). If it produces a fingerprint seen earlier, the loop stops (``Oscillation detected''). Tracking the testbench as well as the RTL matters: a fix that only repairs a wrong expected value in the testbench is a real change.

\begin{viva}
\emph{``What stops the AI from cheating?''} Several things: the testbench is written from the specification, and its prompt tells it the RTL is withheld; a fix may rewrite the testbench freely only when it did not compile, and otherwise may not reduce its checks; if the testbench changed, the original one is run again at the end and a fresh second-opinion testbench is asked for, and if either fails the verdict reads ``PASS, disputed''; with a separate testbench model, a failing check is never rewritten by the design model; the design may not print or end the simulation; and only the simulator's own counts decide the verdict. Also, a fix that changes nothing or repeats an earlier state stops the loop.
\end{viva}

\section{Post-processing}

After the loop, in this order, the orchestrator:
\begin{enumerate}
  \item adds the coverage summary (\code{compute_coverage}), the parsed waveform (\code{parse_vcd}) and the port-level schematic (\code{build_schematic});
  \item runs generic Yosys synthesis if Yosys is installed (if not, the \code{synthesis} event says so; nothing is invented); this also gives the RTL-cell and gate-level netlists for the Schematic tab (Section~\ref{sec:schematicnetlists});
  \item if the run passed but the testbench was changed during repair, re-runs the original testbench and asks for a second opinion (Section~\ref{sec:secondopinion});
  \item for a PASS, measures testbench strength (Section~\ref{sec:strength});
  \item for a PASS, when Yosys, OpenSTA and the cell library are all present, runs the timing and area estimate, which ends with the equivalence proof (Chapter~\ref{ch:timing});
  \item joins every warning (repair stopped because the model became unavailable, a port wider than the check task compares, too few checks, weak testbench, original testbench still failing, second opinion disagrees, the two models disagree, an independent testbench repaired by the design model) into one \code{verification_warning} shown on the result, and sets \code{pass_disputed} when the original or the fresh testbench fails a PASS.
\end{enumerate}
Each extra check is wrapped so that if it fails it only disappears from the result; it can never break a run that already has a verdict. Measured toggle coverage is not a separate stage: the simulator computes it every time it parses a result (Section~\ref{sec:toggle}).

\section{Fallback policy: honest degraded mode}
\label{sec:fallback}

All model calls go through \code{_call_llm} inside \code{pipeline_events}:

\begin{enumerate}
  \item If the run is already offline, call the per-run \code{OfflineProvider}.
  \item Otherwise call the live provider. If it raises an error, \code{_is_upstream_error} decides whether this was the provider's fault (connection error, timeout, rate limit, HTTP 401/403/429/5xx, empty answer). Programming errors are \emph{not} hidden; they surface as errors.
  \item For a provider failure \emph{before any live call has succeeded}, switch this run to offline mode and set a visible \code{fallback_notice}.
  \item For a provider failure \emph{after} a live call succeeded (\code{live_ok}), never switch to the demo. At the repair step the run is kept with the last result (see above); at any other step it fails honestly: ``LLM became unavailable mid-run''. Swapping to a curated design at that point would silently replace your design with an unrelated counter. A specification you confirmed counts as user work, so it sets \code{live_ok} too.
  \item The decision to run offline is made with the provider that serves \emph{the chosen model}, not the default one. Before the last review, a Cloudflare model you picked was silently served by the scripted demo whenever the OpenCode app (the default provider) was missing from the server.
  \item If \emph{you} picked a model in the Model menu (not the server's default) and it fails, the run is never swapped for the scripted demo. \code{_chosen_model_problem} says in plain words why: ``The model you chose (\ldots) could not be used: its usage limit has been reached / it did not answer in time / the provider did not accept the credentials \ldots{} Pick another model from the Model menu.'' The same happens when the separate testbench model fails (``The testbench model you chose \ldots'').
\end{enumerate}

Each run gets its own offline provider instance, so one user's ``buggy demo'' flag cannot leak into another user's run. The helper \code{_invoke} also lets older provider stubs work by dropping the optional \code{reasoning_effort} and \code{spec} arguments only when the function says it does not accept them.

For a demo, this gives a three-step fallback ladder: a live model; then the offline scripted demo (labelled ``Offline demo'' or ``Degraded mode'' on screen); then saved live results opened from the Projects menu. The preflight tool (Chapter~\ref{ch:testing}) checks that all three are ready.

% ---------------------------------------------------------------------------
\chapter{Measuring How Much a PASS Is Worth}
\label{ch:trust}

A PASS only says that the testbench found nothing wrong. It does not say whether the testbench \emph{could} have found something. If the model writes a weak testbench (few checks, wrong expected values, or checks that a repair quietly relaxed), a wrong design can still pass. This chapter describes the five things SiliconScribe does around the loop to show how much a particular PASS is worth.

\begin{analogy}
A student passes a test. Before you trust the mark, you ask: was the test hard enough (strength)? Did it cover the whole syllabus (coverage)? If the teacher changed the answer key halfway through, does a second, fresh paper agree (second opinion)? Was the paper set by someone other than the student (separate testbench model)? Was the question itself what you meant to ask (specification review)?
\end{analogy}

\section{Testbench strength (mutation testing)}
\label{sec:strength}

Mutation testing (Section~\ref{sec:mutation}) is normally run once, by hand, on known designs. \file{backend/mutation.py} runs it automatically on \emph{every} design that passes, using that design's own testbench.

\begin{enumerate}
  \item \code{mask_code} makes a same-length copy of the RTL with comments, strings, compiler directives and parameter declarations blanked out, so nothing in them is ever changed.
  \item \code{candidate_mutants} lists every one-token change the operators can make: an arithmetic, logic or comparison operator swapped (for example \code{+} to \code{-}, \code{&} to \code{|}, \code{==} to \code{!=}, \code{<} to \code{<=}), a \code{~} or \code{!} removed, \code{posedge} swapped with \code{negedge}, a sized constant changed by one, the lowest bit of a binary or hex literal flipped, a 1-bit constant flipped, or an \code{if} condition negated. Bit ranges such as \code{[WIDTH-1:0]} and loop headers are protected: changing them alters a declaration or can create an endless loop, not a realistic bug.
  \item \code{select_mutants} picks up to 12 of them (\code{STRENGTH_MAX_MUTANTS}), spread across the kinds of change and along the file, in a fixed (repeatable) way.
  \item \code{run_strength} first runs the \emph{unmutated} design with the same testbench and times it (the baseline). Then it simulates the same testbench against each mutant, with 6 seconds per mutant and a 40-second budget overall (\code{STRENGTH_BUDGET_SECONDS}). A mutant that does not compile says nothing and is not counted. A mutant the testbench \emph{reports as failing} is \textbf{caught}; one that still passes \textbf{survived}. A mutant that times out counts as caught only if the baseline finished in well under the limit (the mutant made the design hang); otherwise, like a crash, it is \textbf{inconclusive} and not counted.
\end{enumerate}

The result is shown as ``Caught N of M injected bugs'', with each survivor's line number and change. If the time budget runs out, the score is marked as stopped early.

\begin{warnbox}
A survivor is ``not noticed'', not ``proven missed'': some mutants are \textbf{equivalent} (the changed design behaves exactly like the original), so no testbench could catch them. The sample is small and syntactic. The score is evidence about the testbench, not a proof. When at least 5 mutants ran and fewer than half were caught, \code{weak_strength_note} adds a ``Weak testbench'' warning to the PASS; for a small design with fewer than 5 mutants it warns when \emph{none} was caught. Before the last review, a slow testbench with a single check scored 3 of 3 because every mutant ran out of time, and a small design that caught 0 of 3 got no warning at all.
\end{warnbox}

\section{Measured toggle coverage}
\label{sec:toggle}

Coverage asks how much of the design the tests exercised. SiliconScribe measures one kind directly: \textbf{toggle coverage}. A bit is covered when the simulator recorded it going from 0 to 1 \emph{and} from 1 to 0. A bit that never moved was never really tested.

\code{toggle_coverage} in \file{backend/vcd_parser.py} reads the VCD file the simulator wrote (not anything the testbench prints). It counts only signals inside the \emph{design's own} instance: the simulator first finds, in the testbench text, the instance names of modules defined in the RTL (\code{dut_instance_names}), and only signals below one of those instances count, also when the design sits inside a wrapper module or a \code{generate} block. Testbench variables and testbench-side helper modules cannot inflate the number (before the last review, an unrelated stimulus generator in the testbench raised an adder's coverage from 46\% to 77\%). Each run deletes the previous run's waveform first, so an old file can never be measured for a new run. A change through \code{x} or \code{z} does not count. The result is a percentage of bits, the number of bits and signals, and up to 12 signals that never switched (with the bit positions that never rose or never fell). If the VCD was larger than the 8 MB read limit, the figure is marked as a lower bound.

\begin{warnbox}
Only toggle coverage is measured. Line and branch coverage are not. Verilator was tried for them and rejected: with the pip-installed Verilator only 3 of 12 evaluated designs built, and the testbenches of two designs (c3, c4) stopped with ``unsupported tristate construct''. The older \code{compute_coverage} summary (vectors, pass rate, any \code{COV:} lines a testbench prints) is still there, but it only repeats what the testbench reported.
\end{warnbox}

\section{The second opinion}
\label{sec:secondopinion}

The weakest kind of PASS is one that only appeared after a repair \emph{rewrote the testbench}. Maybe the repair fixed a wrong expected value (good). Maybe it relaxed a correct check (bad). The simulator cannot tell which.

So, when the final status is PASS and the testbench is no longer the one first generated, the orchestrator asks for a \textbf{fresh testbench} from the same specification (by the same testbench model), checks it with \code{_tb_problems}, and runs the final design against it. There are three outcomes:
\begin{description}
  \item[Agrees] The fresh testbench passes: ``Second opinion agrees: a fresh testbench, written from the same specification, passes all N of N checks.''
  \item[Disagrees] The fresh testbench fails: the PASS becomes ``PASS, disputed'' and a warning names the failing count.
  \item[Unavailable] The fresh testbench could not be generated, was unusable, did not run cleanly, or ran fewer than 20 checks (too few to count as agreement; a one-check testbench used to count). The reason is shown; nothing is assumed.
\end{description}
The second opinion is skipped in offline mode (the scripted demo has nothing new to say) and can be turned off with \code{SECOND_OPINION_ENABLED=0}.

A real example (s2 prompt, an 8-bit loadable up/down counter, model hy4-flash-free): PASS 8672/8672 after one repair that rewrote the testbench. The fresh testbench also passed (8517/8517), while the original testbench still failed, which was shown as a warning. A hand-written reference then gave 0 mismatches in 40,027 cycles: the design was correct and the original testbench was wrong.

\section{A separate model for the testbench}
\label{sec:independenttb}

By default one model writes both the design and the testbench. Withholding the RTL from the testbench step removes the most direct circular check, but one risk remains: if the model misreads the specification, it can misread it the same way twice, and the wrong design and the wrong testbench agree.

The ``Testbench written by'' menu lets a \emph{different} model, even from a different provider, write the testbench (\code{testbench_model} in the request). Then three rules change:
\begin{itemize}
  \item The design model is \textbf{never} allowed to change the testbench to make a failing check go away: a failing check is reported, not rewritten. Only a testbench whose oracle is broken (it did not compile, or printed no readable summary) may be repaired, and then the result is no longer labelled independent. Two spellings of the same model (\code{x-free} and \code{opencode/x-free}) are not counted as independence.
  \item If the two models disagree (the run ends in FAIL), the result says so and names both models: ``Either the design or the testbench misreads the specification: compare the first failing check with the specification before changing either.''
  \item The result records both models (``Design written by \ldots. Testbench written by \ldots, a separate model''); with one model it says ``Design written by \ldots, which also wrote the testbench''.
\end{itemize}

\begin{keyidea}
Independence makes failures more honest, not passes more frequent. A real example: design by opencode/hy4-flash-free, testbench by nim/poolside/laguna-xs-2.1, 4-bit counter: FAIL 124/146. The counter was correct; the testbench sampled one cycle early. The app refused to pass it and reported that the two models disagree. With one model, a repair could have quietly ``fixed'' either side.
\end{keyidea}

It costs one more model call per run, and with the free models tried it was not reliable enough to be the default, which is why it is an option.

\section{Reviewing the specification}
\label{sec:specreview}

Everything is built from the specification, and the specification is written by the model. If it is wrong (wrong reset polarity, a missing flag), both the design and the testbench will follow it. ``Review the specification before generating'' puts you in the loop at that point:
\begin{enumerate}
  \item \code{POST /api/spec/parse} runs only \code{parse_intent} and returns the specification.
  \item The SpecReview panel shows the module name, the ports table and parameters, and lets you edit the behaviour text and the constraints (one per line).
  \item ``Generate with this specification'' sends it as \code{spec} in the run request. The orchestrator then skips \code{parse_intent} and uses exactly that specification.
\end{enumerate}

A specification from the browser is untrusted input: its names end up in file patterns and prompts. \code{check_user_spec} therefore requires plain identifiers for the module, port and parameter names, 1 to 64 ports, at most 16 parameters, at most 4000 characters of behaviour and at most 30 constraints of 500 characters. The specification actually used is shown read-only on the Results tab under ``Specification the design was built from''.

% ---------------------------------------------------------------------------
\chapter{The Timing and Area Estimate}
\label{ch:timing}

Simulation answers ``does the design do the right thing?''. A chip designer also asks ``how big is it?'' and ``how fast can it run?''. After a PASS, \file{backend/timing.py} gives an early answer to both, using two free tools and a real cell library.

\section{The idea in plain words}

\begin{description}
  \item[Standard cells.] A chip is built from a library of small, pre-designed blocks called \textbf{standard cells}: a NAND gate, an inverter, a flip-flop, and so on. A foundry publishes, for each cell, its area in square micrometres and how long a signal takes to pass through it. That data lives in a \textbf{liberty file} (\file{.lib}).
  \item[sky130.] SiliconScribe uses the open \textbf{SkyWater 130 nm} library, high-density variant (\code{sky130_fd_sc_hd}), at its \textbf{typical corner}: 25\,\textdegree C and 1.8 V. The file is \file{backend/pdk/sky130hd_tt.lib}.
  \item[Mapping (Yosys).] \textbf{Yosys} turns the RTL into a \textbf{netlist}: a list of real sky130 cells and the wires between them. Adding up the cells' areas gives the \textbf{mapped area}.
  \item[Timing (OpenSTA).] \textbf{OpenSTA} is a \textbf{static timing analysis} tool. It does not simulate. It adds up the cell delays along every path from a flip-flop (or input) to the next flip-flop (or output) and compares the longest one, the \textbf{critical path}, with the clock period. The margin left over is the \textbf{slack}: positive means the logic is fast enough for the target clock, negative means it is too slow.
\end{description}

\begin{analogy}
A clock period is the time between two drum beats. Every signal must finish its journey through the gates before the next beat. Slack is how many nanoseconds early the slowest signal arrives. A slack of +9.27 ns at 100 MHz (a 10 ns beat) means the slowest path takes only about 0.73 ns.
\end{analogy}

\section{What one run does}

\code{run_timing(rtl, module, work_dir, target_mhz, testbench_code)} never raises; every outcome is a \code{TimingReport}.
\begin{enumerate}
  \item \textbf{Check the tools}: Yosys and OpenSTA on the PATH and the liberty file present (\code{missing_requirements}). If anything is missing, the run simply has no timing section.
  \item \textbf{Check the input}: a legal module name, the same sandbox scan as simulation, and at most 200 KB of RTL.
  \item \textbf{Apply parameter overrides.} If the testbench instantiates the design as \code{mux #(.WIDTH(16)) dut} or \code{alu #(8) dut}, \code{parameter_overrides} reads those plain numbers (positional ones are matched to the parameters the design declares) so the circuit timed is the circuit that was simulated. Expressions are never guessed; they are reported as ``not applied''.
  \item \textbf{Map with Yosys} (\code{_map_with_yosys}): read the design, \code{proc; tribuf} (so a tri-state survives as a cell that can be detected), \code{synth -flatten} (hierarchy is flattened, so area covers every sub-module), map flip-flops and logic to sky130 cells (low-power isolation cells excluded), write the statistics and \file{netlist.v}.
  \item \textbf{Read what was really built} (\code{analyze_netlist}): which input ports drive the flip-flops' clock pins, whether any unmapped cell, latch or clock-gating cell is present.
  \item \textbf{Refuse instead of guessing} (next section).
  \item \textbf{Time with OpenSTA}: write a small script that reads the library and netlist, creates one clock per clock port found in the netlist at the target period, sets input and output delays to zero, and reports the worst setup path. A design with no flip-flops is timed against a virtual clock, so the answer is its longest input-to-output path.
  \item \textbf{Parse the report} (\code{parse_mapped_stats}, \code{parse_sta}): area, sequential area, cell count, flip-flops, the most common cell types, worst slack, maximum frequency ($1000 / (\text{period} - \text{slack})$ MHz), the critical path's start, end and delay, and its pin-by-pin rows (\code{critical_path_rows}, Section~\ref{sec:critpath}).
  \item \textbf{Prove the netlist equals the RTL} (\code{prove_equivalence} in \file{equivalence.py}, Section~\ref{sec:equivalence}). This runs unless \code{EQUIVALENCE_ENABLED=0}. If the proof itself crashes, the timing numbers stay and the line says the check could not be run.
\end{enumerate}

\section{Clocks come from the netlist}

It would be easy to guess the clock from a port name such as \code{clk}. \file{timing.py} does not. It reads the mapped netlist and takes whatever input port is actually wired to the flip-flops' clock pins, whatever its name. A design with several clocks gets one clock each. Every clock's own paths are timed, \emph{including} the paths from the inputs and to the outputs of that clock's flip-flops; only paths that cross from one clock to another are not timed, and the note says so.

\begin{warnbox}[A bug the last review found here]
Input and output delays used to be attached only to the first clock. Paths from the inputs into the second clock's flip-flops were therefore not timed at all, so a slow multiplier on a second clock was reported as meeting 500 MHz (slack $+1.74$ ns) when on its own it misses by 1.43 ns. Now every clock gets its own input and output delays.
\end{warnbox}

\section{When the estimate refuses}

Some designs cannot be timed honestly with this simple setup. For them the report gives a reason instead of a number:
\begin{itemize}
  \item \textbf{tri-state buffers}: Yosys would otherwise quietly turn them into plain wires, so the circuit timed would not be the one simulated;
  \item \textbf{latches}: latch-based timing is not estimated;
  \item \textbf{combinational loops}: path timing is not defined (OpenSTA reports a disabled loop edge);
  \item \textbf{derived or gated clocks}: a flip-flop clocked by internal logic or a clock-gating cell rather than an input port;
  \item \textbf{no logic cells}: a design that maps to only wires or constants has nothing to time;
  \item \textbf{not synthesizable as written}: Yosys could not map it.
\end{itemize}

\section{What the numbers are, and are not}

\begin{warnbox}[Pre-layout estimate: not sign-off]
The app prints this with every result: one typical corner; no placement or routing, so \textbf{no wire delay}; an ideal clock (no clock tree); zero input and output delays; \textbf{setup analysis only} (hold needs a real clock tree to mean anything). \textbf{Power is not reported.} Real chips are checked after place and route, at several corners. So this answers ``is this logic anywhere near the target clock, and roughly how big is it?'', nothing more. Never present it as final (sign-off) timing or as the area of a manufactured chip.
\end{warnbox}

Two numbers need care. The maximum frequency is computed as $1000 / (\text{period} - \text{slack})$, which is only right when the worst path takes one whole clock period. When the worst path starts or ends on the falling clock edge (a half-cycle path), that formula gives a number that changes with the target clock, so no maximum frequency is shown at all. And a design without a clock is timed against a virtual clock: its longest input-to-output delay is a real number, but $1/\text{delay}$ is not a clock rate, so the chat message does not quote one.

The cell counts come from Yosys' statistics. Yosys writes large areas in exponent form (\code{1.6E+03} for a group of 64 flip-flops), and the parser used to skip such lines, so a 64-bit register was refused as ``no logic cells'' and larger counters reported 0 flip-flops; it now reads exponent numbers. If the testbench sets parameters with \code{defparam}, the note says that those values were \emph{not} applied, so the numbers are for the design's default parameters.

Examples from the built-in designs: the 4-bit counter maps to 11 cells, 150.144\,\umsq{}, minimum period 0.727 ns (meets 100 MHz with +9.27 ns slack); the 8-bit ALU has no clock, 109 cells, 636.861\,\umsq{}, longest path 2.69 ns. The live-generated FIFO (s6) is 120 cells, 1830.51\,\umsq{}, minimum period 1.845 ns.

\section{The critical path, pin by pin}
\label{sec:critpath}

OpenSTA's report lists the worst path as a column of rows, one per pin the signal passes. \code{parse_sta} picks the path block with the smallest slack, and \code{critical_path_rows} turns its data part into rows of five fields: the \textbf{pin} (for example \code{_127_/X}, output X of cell number 127), the sky130 \textbf{cell} type with the library prefix removed (\code{maj3_1}), whether the signal \textbf{rises or falls} there, the \textbf{incremental delay} added by that stage, and the \textbf{arrival} time so far. Rows after ``data arrival time'' (the clock side of the check) are not part of the data path and are dropped; at most 60 rows are kept. The rows travel to the browser in \code{TimingReport.critical_path}.

On the built-in 8-bit ALU (no clock) the worst path starts at input \code{b[0]}, goes through one \code{nor2b_1}, then six \code{maj3_1} cells in a row (a 3-input majority gate is exactly the carry logic of a full adder, so this is the ripple-carry chain, about 0.36 to 0.39 ns per stage), then \code{o21ai_0} and \code{o22ai_1}, and ends at the output \code{overflow} after 2.69 ns. Section~\ref{sec:readtiming} shows how to read this in the app.

\section{Formal equivalence: proving the netlist is the RTL}
\label{sec:equivalence}

The timing and area numbers describe \file{netlist.v}, the mapped netlist. They are only worth something if that netlist is the design that was simulated. \file{backend/equivalence.py} checks this with a formal proof (Section~\ref{sec:formalbasics} explains the ideas). It uses only Yosys and its built-in SAT solver; no external solver is installed.

\code{prove_equivalence(work_dir, module, liberty, params)} runs in the timing stage's own folder, which already holds \file{design.sv} and \file{netlist.v}:
\begin{enumerate}
  \item \textbf{Read the RTL as ``gold''}: \code{read_verilog -sv design.sv}, the \emph{same parameter overrides} the netlist was mapped with (\code{hierarchy -chparam}), \code{proc}, \code{flatten}, and \code{memory}. The \code{memory} pass turns arrays and lookup tables into plain logic and registers, as synthesis does, so such designs are proven instead of skipped.
  \item \textbf{Read the netlist as ``gate''}: first \code{read_liberty} on the sky130 library, so every cell's behaviour comes from the logic function written in the liberty file itself, not from a model someone typed by hand. Then \code{read_verilog netlist.v} and \code{flatten}.
  \item \textbf{Decide the method}: if the netlist contains a sky130 flip-flop cell, the design is sequential; otherwise it is combinational.
  \item \textbf{Combinational: miter and SAT.} \code{miter -equiv} builds one circuit that feeds the same inputs to gold and gate and raises a flag if any output differs. \code{sat -prove-asserts -set-def-inputs -enable_undef} asks the solver whether any input can raise the flag. Inputs are limited to real 0 and 1 values (\code{x} is not a value a real input pin carries), and outputs where the RTL itself says ``don't care'' (\code{x}) are not compared. ``SUCCESS'' means proven over every input value; ``FAIL'' means the solver found inputs where they disagree.
  \item \textbf{Sequential: induction.} \code{equiv_make} pairs the registers of gold and gate by name, \code{async2sync} turns asynchronous resets into a form the prover handles, then \code{equiv_simple -seq 5} and \code{equiv_induct -seq 5} try to prove every pair and every output by temporal induction. \code{equiv_status} reports how many points were proven.
\end{enumerate}

\code{interpret} turns Yosys' output into one of four results, stored as \code{TimingReport.equivalence} with a plain sentence in \code{equivalence_method}:
\begin{description}
  \item[\code{proven}] Combinational: ``SAT proof over every input value''. Sequential: every register and output point proven by induction.
  \item[\code{different}] The solver found input values where the netlist and the RTL disagree. The app shows a red alert and declares the timing and area invalid.
  \item[\code{not_proven}] Induction proved only some points. This usually means synthesis re-encoded the state of a state machine, so registers no longer pair by name. It is also the result, with the reason, when a netlist flip-flop uses a clock edge the RTL does not use, when there are several clocks, or when a register gets a starting value from its declaration or an \code{initial} block (\code{_sequential_limits}): induction would otherwise ``prove'' those wrongly. It is never shown as a pass.
  \item[no result] Yosys could not run the check, the design has more than 600 mapped cells (the proof is skipped: SAT on a large multiplier can run for minutes), or the proof took longer than 20 seconds and was stopped. The app shows ``Netlist not formally checked'' with the reason, never a silent gap.
\end{description}

\textbf{Measured result.} On the 14 evaluated designs plus 4 built-in designs (18 in all): 16 \textbf{proven}, 2 \textbf{not proven} (f1 and f2, the two state machines, whose state was re-encoded by synthesis) and 0 different. A deliberately sabotaged netlist, with one NAND gate swapped for a NOR, was reported as different. The proof shows netlist $=$ RTL for those 16 designs; it says nothing about whether their RTL meets the prompt.

\begin{warnbox}[What the proof does not cover]
It does not prove that the RTL is correct; a bug in the RTL is copied faithfully into the netlist. ``Not proven'' is an honest ``don't know'', not a hidden failure: for f1 and f2 the gate-level simulation check below still passed with the same testbench. The last review checked the proof's blind spots directly: a netlist with one flip-flop moved to the falling clock edge, and a design whose register started at 5 from its declaration, were both ``proven'' before; both are now ``not proven'' with the reason.
\end{warnbox}

\section{How the estimate itself was checked}

\file{backend/tools/verify_mapping.py} checks the timing stage with three tests that do not read the number they check from the same place:
\begin{enumerate}
  \item \textbf{Gate-level simulation}: Yosys writes simulation models of the sky130 cells straight from the liberty file; the mapped netlist is simulated in Icarus with the \emph{same} testbench as the RTL and must give identical Passed/Failed counts. This shows the netlist computes what the RTL computes.
  \item \textbf{Area recount}: area and cell count are recomputed from \file{netlist.v} and the liberty file and compared with what \code{run_timing} reported.
  \item \textbf{Timing sanity}: the design is timed at a second target; the implied minimum period ($\text{period} - \text{slack}$) must not change.
\end{enumerate}
On 19 designs (the 5 built-in ones and the saved evaluation designs) all 19 gave identical pass counts in gate-level simulation. A design the estimate refuses on purpose is reported as REFUSED with its reason, not as a failure. The checker was also given deliberately broken netlists and caught them.

The two checks complement each other: the simulation covers every design, including the state machines, but only the testbench's inputs; the proof covers every input value, but cannot always finish for re-encoded state.

\section{The tool versions}

Yosys 0.71 runs as a WebAssembly build installed with \code{yowasp-yosys}, reached through a small \code{yosys} shim in \file{~/.local/bin}. OpenSTA 3.2.0 was built from source (\file{backend/tools/build_opensta.sh}) and is reached as \code{sta}. The orchestrator runs timing only when \code{TIMING_ENABLED} is on (the default) and \code{missing_requirements()} is empty. \code{POST /api/timing/run} runs the same function on its own.

% ---------------------------------------------------------------------------
\chapter{The LLM Service}
\label{ch:llm}

\file{backend/llm_service.py} talks to the AI providers; \file{backend/model_catalog.py} decides which models the Model menu offers. No other file knows about providers.

\section{Choosing the default provider}

At start-up, \code{_make_provider} reads environment variables (from \file{backend/.env}, which wins over stale shell variables) and picks, in order:

\begin{center}
\small
\begin{tabularx}{\textwidth}{@{}clX@{}}
\toprule
\textbf{Priority} & \textbf{Provider} & \textbf{Selected when} \\
\midrule
0 & Offline & \code{OFFLINE_MODE=1} (always wins, even if keys are set). \\
1 & OpenCode app & \code{LLM_PROVIDER=opencode-app} and the \code{opencode} command exists. Model from \code{LLM_MODEL}. This is the set-up used for the live demo and evaluation. \\
2 & OpenCode Zen & \code{OPENCODE_API_KEY} or \code{ZEN_API_KEY} is set. Model from \code{ZEN_MODEL} or \code{LLM_MODEL}. \\
3 & NVIDIA (older single-model path) & \code{NVIDIA_API_KEY} is set. \\
4 & OpenAI or compatible & \code{OPENAI_API_KEY} is set; add \code{OPENAI_BASE_URL} for OpenRouter, Groq, a local server, etc. \\
5 & Offline & Nothing configured. \\
\bottomrule
\end{tabularx}
\end{center}

Placeholder values such as \code{your_key_here} are treated as empty. Rows 2--4 are an \code{OpenAIProvider}, configured with a base URL and model name.

\section{The OpenCode app provider}

OpenCode is a coding assistant app that the user installs and signs in to with their own account; its free models are meant to be used from inside the app. \code{OpenCodeAppProvider._chat} therefore drives the app's own scripting command instead of imitating its API:
\begin{lstlisting}[style=plain]
opencode run -m <model> --agent plan --archived --format json "<prompt>"
\end{lstlisting}
Every call runs in a fresh, empty temporary folder with the read-only \code{plan} agent (change it with \code{OPENCODE_AGENT}), and with an environment that has the Cloudflare and NIM credentials removed (\code{model_catalog.child_env}). The JSON lines the app prints are read; the text parts become the answer, an error event becomes an upstream error. It is a local, single-user convenience: not something to expose on a public server.

\section{Model routing: one menu, three sources}
\label{sec:routing}

Each run can use a different model, chosen in the Model menu. The model id itself says which provider serves it; \code{get_provider(model)} routes it:

\begin{center}
\small
\begin{tabularx}{\textwidth}{@{}lp{3.7cm}X@{}}
\toprule
\textbf{Id looks like} & \textbf{Served by} & \textbf{Notes} \\
\midrule
\code{opencode/deepseek-v4.1-flash-free} & The configured provider (normally the OpenCode app) & Every model the installed app lists with \code{opencode models}, about 460. A model counts as free only when its name ends in \code{-free} or \code{:free}; the rest are hidden behind ``Show other models'' and are never claimed to be paid, because the price cannot be known from here. \\
\code{github-copilot/auto} & The OpenCode app, after Copilot is connected in it & On the Free and Student plans only ``auto'' exists: GitHub picks the model and does not say which, so the result cannot name its model and runs are not reproducible. It uses your Copilot allowance. \\
\code{@cf/openai/gpt-oss-120b} & \code{CloudflareProvider} (Cloudflare Workers AI) & Live catalogue of text-generation models on your account (40 when checked). Needs \code{CLOUDFLARE_ACCOUNT_ID} and \code{CLOUDFLARE_API_TOKEN}. \\
\code{nim/poolside/laguna-xs-2.1} & \code{NimProvider} (NVIDIA NIM hosted API) & Live catalogue of chat models (15 when checked); the \code{nim/} prefix keeps them apart from OpenCode's own \code{nvidia/...} ids and is removed before the request. Needs \code{NVIDIA_NIM_API_KEY} (a key starting \code{nvapi-}); no other NVIDIA variable is read for this. \\
\bottomrule
\end{tabularx}
\end{center}

Cloudflare and NIM both use the same OpenAI-compatible chat code as every other provider. They get no \code{reasoning} request field (both are strict about unknown fields) and no per-stage model routing (the \code{LLM_MODEL_<STAGE>} variables name OpenCode models). If one of them is chosen but not set up, the run stops with ``Cloudflare is not set up\ldots'' or ``NVIDIA NIM is not set up\ldots''.

\begin{keyidea}
Keys live only in \file{backend/.env} (git-ignored) or the environment of the terminal that starts the backend. The web page never asks for, stores or receives a key. \code{cloudflare_status} and \code{nim_status} return models and hints, never the key. Simulation, synthesis and timing tools get a scrubbed environment, and the OpenCode app gets an environment without the Cloudflare and NIM credentials.
\end{keyidea}

\section{The model catalogue (\texttt{model\_catalog.py})}

\begin{description}
  \item[\code{opencode_models}] Runs \code{opencode models} in an empty temporary folder (40 s limit) and parses the \code{provider/model} lines with \code{parse_opencode_models}: free ones first within each provider, GitHub Copilot entries carry the note above. A good listing is cached for 10 minutes, a failed one for 30 seconds.
  \item[\code{cloudflare_status}] Reads up to 5 pages of Cloudflare's model search, keeps text-generation models, and caches them for 30 minutes. If the list cannot be read it shows a short suggested list with a hint saying so.
  \item[\code{nim_status}] Reads NVIDIA's \code{/models} list and drops entries that are not chat models (embeddings, speech, safety classifiers, image models) and one known to answer 404. Models that answered quickly in a live test are listed first.
  \item[\code{list_models}] (in \file{llm_service.py}) Combines the three: the OpenCode list (free ones only unless \code{include_paid}), then Cloudflare, then NIM, plus a \code{providers} list saying which tabs are available and why not. The configured default model is always first. \code{GET /api/models?refresh=true} re-reads all catalogues and is rate limited, because it calls out.
\end{description}

Model names change often (NVIDIA retired DeepSeek V4.1 Flash on 16 Nov 2026 with HTTP 410), which is why the menu reads the catalogues live instead of a hand-written list.

\section{Checking model ids}
\label{sec:modelid}

A model id from the browser ends up as an argument of a subprocess (\code{opencode run -m <id>}) or part of a URL. An id starting with \code{-} could be read as a command-line flag. So every id must match
\begin{lstlisting}[style=plain]
^[A-Za-z0-9@][A-Za-z0-9._:/@+-]{0,99}$
\end{lstlisting}
(no leading dash, no spaces or shell characters, at most 100 characters). The check is done twice: by Pydantic on every request (\code{model} and \code{testbench_model}) and again by \code{is_valid_model_id} just before the subprocess is started.

\section{One chat call: \texttt{OpenAIProvider.\_chat}}

\begin{enumerate}
  \item \textbf{Pick the model.} A model the user chose in the UI (if different from the default) wins; otherwise the stage-specific model from \code{LLM_MODEL_INTENT}, \code{_RTL}, \code{_TESTBENCH}, \code{_EXPLAIN} or \code{_FIX}; otherwise the default (or the cheaper ``mini'' model for the explanation).
  \item \textbf{Control hidden reasoning.} If no reasoning effort was requested, send \texttt{reasoning: \{enabled: false\}}, because hidden reasoning tokens count against the answer budget. Gateways that reject this field get one retry without it.
  \item \textbf{Send} with the configured timeout (default 120 s, \code{LLM_TIMEOUT_SECONDS}) and token budget (default 8192, \code{LLM_MAX_TOKENS}).
  \item \textbf{Classify failures.} Provider errors are re-raised with the prefix \code{LLM_UPSTREAM:} so the orchestrator can recognise them.
  \item \textbf{Reject bad answers.} If the model stopped because it hit the token limit (\code{finish_reason = "length"}) or returned empty content, raise an error. The code deliberately never uses the model's hidden reasoning text as an answer; that once produced a half-written function shipped as a testbench.
\end{enumerate}

\section{The five stage functions}

\begin{center}
\small
\begin{tabularx}{\textwidth}{@{}lcX@{}}
\toprule
\textbf{Function} & \textbf{Temp.} & \textbf{What it sends and returns} \\
\midrule
\code{parse_intent} & 0.3 & The prompt; returns JSON validated as \code{RTLDesignSpec}. \\
\code{generate_rtl} & 0.2 & Module name, parameters, ports, behaviour, constraints, optional target frequency; returns the extracted module. \\
\code{generate_testbench} & 0.3 & The same specification, \emph{without the RTL}, with an explicit note that the RTL is withheld; returns the extracted testbench. \\
\code{explain_design} & 0.5 & The RTL; returns 2--3 sentences. \\
\code{fix_design} & 0.2 & Spec, RTL, testbench, last 4000 log characters; returns new code, summary, diagnosis, \code{fix_type = "llm_diagnosis"}. \\
\bottomrule
\end{tabularx}
\end{center}

The full system prompts are in \file{backend/llm_service.py} (and Appendix~B of the project report). The testbench prompt insists on Icarus-compatible Verilog-2001, a top module named \code{testbench}, a reset sequence first, directed corner cases plus exhaustive testing when there are at most 4096 input combinations (otherwise at least 100 random vectors), every comparison through the standard \code{chk} task (which uses \code{===} and prints one \code{CHK|...} line per check), the exact \code{Passed:}/\code{Failed:}/\code{ALL TESTS PASSED} summary lines, and a watchdog that prints a failure and stops a hung run.

\section{Reading the model's answer}
\label{sec:llmparse}

Models wrap code in prose and Markdown fences, think out loud first, and sometimes give a draft before the real answer. The parsers are forgiving, but in a careful order:

\begin{description}
  \item[\code{_strip_think}] Removes reasoning blocks (\code{<think>...</think>}), and everything before a lone closing \code{</think>} (some models print only the closing tag). Every parser below starts with this, so a draft written while ``thinking'' is never taken as the answer.
  \item[\code{_strip_code_fence}] Finds every complete \code{module ... endmodule} unit in the whole reply and keeps \emph{one definition per module name: the last one} (\code{_last_definition_of_each}). A reply with a draft and then a ``corrected'' copy used to produce the module twice, which Icarus rejects as ``already declared''; helper modules with other names are all kept.
  \item[\code{_extract_json}] Takes the \emph{last} complete top-level JSON object in the reply that parses, accepting raw newlines inside strings (models often paste Verilog into JSON without escaping). The old ``first brace to last brace'' slice broke on prose such as ``\{as requested\}'' before the JSON, or on a draft object followed by the final one.
  \item[\code{_parse_fix_response}] Reads the \emph{last} \code{<diagnosis>}, \code{<rtl>}, \code{<testbench>} and \code{<summary>} sections (a model that first restates the format inside its reasoning no longer fools it); falls back to JSON, then to \code{_salvage_fix_json}, which recovers each field by locating the next known key. If none of these yields any code, the answer is marked \code{parse_failed} (Section~\ref{ch:pipeline} describes the one retry).
\end{description}

The specification model is also forgiving about spelling: port directions such as \code{Output}, \code{out} or \code{in} are normalised, a missing description becomes empty text, a behaviour given as a list is joined, and a single constraint string becomes a list. A module name written by the model is turned into a plain identifier (\code{model_module_name}: \code{Up/Down Counter} becomes \code{Up_Down_Counter}, and \code{testbench}, which would collide with the testbench's own module, becomes \code{testbench_dut}). A name \emph{you} type in the specification review is not rewritten; it is checked and refused with a clear message if it is not a plain identifier.

When a call fails, the reason shown is read with whole-word patterns (\code{_reason}): HTTP 429 or ``too many requests'' is a usage limit; 410 or ``retired'' means the provider retired the model; an empty answer after reasoning means the model spent its whole answer budget on thinking; a missing \code{opencode} command is reported as such. The earlier substring tests gave wrong reasons, for example ``credentials'' for any message that contained \code{max_tokens}. The HTTP client retries a failed call twice (\code{LLM_MAX_RETRIES}); three retries with a 120-second timeout could make one call wait eight minutes.

\section{The offline provider}

\code{OfflineProvider} needs no internet. It is used when \code{OFFLINE_MODE=1}, when nothing is configured, or as the labelled fallback described in Chapter~\ref{ch:pipeline}. It matches whole words in the prompt (\code{match_design}): ``alu'' with ``8-bit'' gives the 8-bit ALU, then ``alu'', ``counter'', ``adder'' or ``mux'' (plurals allowed); anything else defaults to the counter, with a notice saying the request did not match. Whole words matter: plain substring matching sent ``8-bit counter with a load \emph{value}'' to the ALU (``value'' contains ``alu'') and ``dis\emph{count} calculator'' to the counter, both labelled as exact matches. A ``half adder'' is answered with the built-in full adder but labelled as not an exact match. A prompt containing ``buggy'' gets the counter with a deliberate bug (counting by two), which demonstrates the repair loop. Its ``fix'' simply swaps in the known-correct design, so it reports \code{fix_type = "offline_scripted"} and \emph{no} diagnosis, because inventing a root cause would be dishonest.

\section{The model list for other providers}

With the OpenCode app, the menu lists what the app offers (Section~\ref{sec:routing}). With any other configured provider (Zen, NVIDIA through \code{NVIDIA_API_KEY}, OpenAI-compatible), \code{list_models} offers a short curated list, or only the configured model if it is not in that list, so the UI does not send a model name the provider cannot serve.

% ---------------------------------------------------------------------------
\chapter{The Simulator and Sandbox}
\label{ch:simulator}

\file{backend/simulator.py} runs untrusted code, so it is written defensively.

\section{The path of one simulation}

\code{IcarusSimulator.simulate(design_id, rtl, tb, timeout)}:

\begin{enumerate}
  \item \textbf{Validate the design id}: 1--32 characters of letters, digits, \code{_} or \code{-}, and the resolved folder must stay inside the workspace. This blocks tricks such as \code{../../etc}.
  \item \textbf{Policy scan} (\code{sandbox_violations} plus \code{rtl_integrity_violations}); any violation returns status ERROR with the reasons, before any file is written.
  \item \textbf{Wait for the design's folder, then for a slot}: only one run at a time may use a design's folder (\code{_dir_lock}), and at most 2 simulations run at once (\code{SIM_MAX_CONCURRENT}); extra requests wait their turn. Without the folder lock, two runs of the same design (two tabs, a double-clicked Re-run) overwrote each other's files and got each other's verdicts: the last review measured a wrong verdict in 25 of 25 tries.
  \item \textbf{Remove the previous run's outputs}: old waveform files, transcript and compiled simulation are deleted, so an old file can never be reported for this run (a testbench that dumps to its own file name used to leave the old \file{design.vcd} behind, and it was shown as the new waveform).
  \item \textbf{Write files} (\code{write_files}): each file at most 200 KB; \code{design.sv} and \code{testbench.sv} in the design's folder; \code{ensure_vcd_dump} inserts \code{$dumpfile("design.vcd")} and \code{$dumpvars(0, testbench)} right after the first module header that is not inside a comment, if the testbench has no dump call of its own (a dump call mentioned only in a comment does not count).
  \item \textbf{Compile} (\code{compile}): \code{iverilog -g2012 -Wall}, 30-second limit, scrubbed environment, under the same CPU, memory and file limits as the simulation; a compile that times out is stopped together with its helper process \code{ivl} (\code{run_tool} kills the whole process group). A one-line file with \code{`timescale 1ns/1ps} is compiled first so every file has a sensible default time unit without shifting line numbers. Failure returns ERROR with parsed file/line errors. A port-width mismatch between testbench and design, which Icarus only warns about, is treated as an error, because a narrower wire silently hides the upper bits.
  \item \textbf{Run} (\code{run}): \code{vvp} with the timeout clamped to 1--120 s, an output-size cap (\code{ulimit -f}), a CPU-time cap (\code{ulimit -t}) and a memory cap (\code{ulimit -v}, 1024 MB by default, enforced on Linux only), stdin closed, transcript to \file{transcript.log}. Overrun returns TIMEOUT, with the last output kept so the fixer can see where it hung.
  \item \textbf{Decide the verdict} (\code{parse_results}), read the \code{CHK} rows (\file{checks.py}) in one pass that keeps only counters and the rows that can be shown (a runaway testbench printing millions of lines used to take about 470 MB here; it now takes about 2 MB), and measure toggle coverage of the design's own instance from the VCD.
\end{enumerate}

\section{How the verdict is decided}
\label{sec:verdict}

The testbench prints its own summary, so \code{_decide_verdict} must not trust the first or the loosest match. Each rule below came from a reproduced review finding.

\begin{center}
\begin{tabularx}{\textwidth}{@{}Xl@{}}
\toprule
\textbf{Condition in the transcript} & \textbf{Status} \\
\midrule
No whole-line \code{Passed: <number>} and \code{Failed: <number>} pair, and no failure lines & ERROR \\
No valid summary, but failure lines (FAIL, ERROR, FATAL, \code{CHK|FAIL}) & FAIL \\
\code{Failed: N} with N $>$ 0 anywhere, even if a later line says \code{Failed: 0} & FAIL \\
\code{Failed: 0}, but the log has lines starting with FAIL, ERROR or FATAL, or a \code{CHK|FAIL} row & FAIL \\
\code{Passed: 0}: no checks were executed & ERROR \\
The simulator exited with a non-zero code (killed by a limit, crashed) & ERROR \\
Otherwise (\code{Passed} $>$ 0, \code{Failed: 0}, clean exit) & PASS \\
\bottomrule
\end{tabularx}
\end{center}

So a PASS needs a valid summary \emph{and} at least one executed check, and nothing in the log may say otherwise. Zero checks is never a PASS. The words \code{ALL TESTS PASSED} alone never decide anything. Only ASCII digits count as numbers (a \code{²} once crashed the parser).

Three rows of the table came from the last review. A testbench that resets its counters for each phase could print \code{Failed: 4} and later \code{Failed: 0}, and only the last line counted. A \code{CHK|FAIL} row was ignored next to \code{Failed: 0}. And failure lines without a summary gave ERROR, which the repair loop treats as ``the testbench is broken, rewrite it'' (Section~\ref{ch:pipeline}): a testbench that stopped at its first mismatch with \code{$fatal} could be rewritten away. All three now give FAIL.

Test count is passed plus failed. Up to 20 FAIL/ERROR lines are collected as errors. The \code{log_excerpt} the debugger sees is the first 1000 and the last 1500 characters of the transcript (with passing \code{CHK} lines removed), so the first failing vector is not cut off.

The design under test must not be able to write the verdict either. \code{rtl_integrity_violations} rejects RTL that contains \code{Passed:}, \code{Failed:}, \code{ALL TESTS PASSED}, \code{$finish}, \code{$stop} or \code{$fatal}; that prints anything (\code{$display}, \code{$write}, \code{$strobe}, \code{$monitor}) or has a \code{final} block; or that reaches into the testbench (\code{testbench.}\ldots, also with a comment before the dot, or \code{$root}). A design that printed a fake summary at time 0 was once reported as PASS. The last review found two more ways: a \code{final} block (which runs after \code{$finish}, so its lines come last) printing \texttt{\%cassed: 16} with the character code of P, and a macro alias \code{`define T testbench} used as \code{`T.fail_count = 0}. Both are refused; the macro case because the rules are run again on the macro-expanded RTL. Comments are ignored for these checks, so a comment such as ``// final value'' is fine.

\section{The sandbox policy}
\label{sec:sandbox}

Generated Verilog runs with the server's permissions, and Verilog has file-access system tasks. The final review proved the risk: a planted secret file was read by a testbench and returned in the API response. The protection has these layers:

\begin{enumerate}
  \item \textbf{Static policy.} Reject \code{$fopen}, \code{$fwrite}, \code{$fdisplay}, \code{$fgets}, \code{$fscanf}, \code{$fread}, \code{$readmemh}/\code{b}, \code{$writememh}/\code{b}, \code{$system}, \code{$sdf_annotate}, \code{$dumpports} and their variants; any \code{`include}; macro token pasting; and any \code{$dumpfile} that is not a plain file name in quotes. Forbidden names are searched in the raw text, comments included; checks that span tokens (the \code{$dumpfile} argument, a reference to the testbench) run on text whose comments were removed by a string-aware scanner, so \code{$dumpfile/**/("/abs/x")} is caught too.
  \item \textbf{Allow-list of system tasks.} A list of forbidden tasks always misses one (a review found a testbench that wrote a file with \code{$fopena} and \code{$fputc} and still got PASS). So every system task must also be on an allow-list of side-effect-free tasks: printing, simulation control, time and random numbers, waveform dump, string formatting and maths. Anything else is refused.
  \item \textbf{Checked again after macro expansion.} A macro can build a name the scan never sees: \code{`define ID(a) a} followed by \code{`ID($)fopen} wrote a file in the last review. Whenever a macro is used, Icarus' own preprocessor expands the RTL and testbench together (in compile order) and every rule runs again on the result; code that cannot be preprocessed is refused. Legitimate macros, such as a check macro that calls \code{$display}, still work.
  \item \textbf{The same check in several places.} The scan runs before compiling, before the linter calls its tools, and before the timing stage maps the design, not only before simulation.
  \item \textbf{Scrubbed environment.} \code{clean_env} passes only \code{PATH}, \code{LANG}, \code{LC_ALL} and \code{TMPDIR} to the tools, so API keys never reach them.
  \item \textbf{Resource caps.} The compiler and the simulator run under the same caps, and a tool that times out is stopped together with its child processes. \code{ulimit -f} limits any file the simulation writes (about 64 MB), so an endless \code{$display} loop cannot fill the disk; CPU time and memory are capped (memory on Linux only); at most 2 simulations at once.
\end{enumerate}

\begin{warnbox}
This is defence in depth, not a jail. A static scan of a language as flexible as Verilog cannot be proven complete. For many users on one server, each simulation should run in a container or micro-VM. The report lists this as a limitation.
\end{warnbox}

% ---------------------------------------------------------------------------
\chapter{Supporting Modules}
\label{ch:support}

\section{Linter (\texttt{linter.py})}

\code{VerilogLinter.lint(rtl, tb)} returns \code{ok}, \code{errors}, \code{warnings} and raw output. It tries Verilator (\code{--lint-only -Wall}), then Icarus (compile to \file{/dev/null}), and always runs pattern heuristics (\code{_heuristic_checks}) such as unbalanced \code{module}/\code{endmodule}. With no tool installed, heuristic errors stay errors, so broken code is never reported as clean. The width heuristic compares a number only with the width of the signal it is \emph{assigned} to (\code{count <= 16;} for a 4-bit \code{count}); it used to compare any number with the smallest vector in the file, so \code{if (cnt <= 9)} next to a 2-bit select was flagged. The orchestrator runs lint three times (RTL alone, RTL plus testbench, after each fix) and only reports; lint never blocks the pipeline. Old lint is cleared when a repair changes the code, and when the simulator compiled and ran the code, the Results tab labels lint as \emph{advisory}.

\section{Coverage (\texttt{cov\_metrics.py})}
\label{sec:coverage}

\code{compute_coverage} always returns test vectors, passed, failed, pass rate and \code{bins}. Bins come from lines the testbench prints as \code{COV: <bin> <count>}. If the testbench prints figures such as ``Toggle coverage: 85\%'', they are picked up as well. It is honestly scoped: it summarises what the testbench reported. The \emph{measured} figure is separate: \code{toggle_coverage} in \file{vcd_parser.py} reads the simulator's VCD and is stored as \code{measured_toggle} on every simulation result (Section~\ref{sec:toggle}). The toggle figure on the Results tab is always the measured one.

\section{VCD parser (\texttt{vcd\_parser.py})}
\label{sec:vcd}

\code{parse_vcd} reads at most 8 MB of the VCD, the timescale, and every \code{$var} declaration (stripping \code{[3:0]} from names, since width is stored separately). It then walks the value changes: \texttt{\#t} sets the time, \code{0!}/\code{1!} are single-bit changes, \code{b0101 code} are bus changes (converted to hex by \code{_bin_to_hex}; any \code{x} or \code{z} bit makes the whole value \code{x} or \code{z}). Limits: 40 signals and 2000 changes per signal, with flags telling the UI when something was truncated.

\section{Schematic and synthesis}
\label{sec:schematicnetlists}

\code{build_schematic} groups the specification's ports into inputs, outputs and inouts; the frontend draws a block with ports on each side (every inout included, with room for the longest port name). \code{YosysSynthesizer} runs \code{read_verilog; hierarchy; proc; opt; write_json; synth -flatten; check; stat; write_json} with a 30-second limit, \emph{without} a cell library, using relative file names, a clean environment and closed input. \code{-flatten} matters: without it a ripple adder built from four \code{full_adder} instances was reported as ``4 cells of type full\_adder'', and a latch inside a sub-module was missed. The two \code{write_json} steps save two snapshots of the design for the Schematic tab:
\begin{itemize}
  \item after \code{proc; opt}: \textbf{word-level RTL cells} such as \code{$add}, \code{$sub}, \code{$mux}, comparators and \code{$adff} flip-flops (field \code{rtl_netlist});
  \item after \code{synth}: \textbf{single-bit generic gates} such as \code{$_XOR_}, \code{$_NAND_} and \code{$_DFFE_PN0P_} (an enabled flip-flop with an active-low asynchronous reset) (field \code{json_netlist}). These are Yosys' own generic cells, not sky130 standard cells; the sky130 netlist exists only inside the timing stage.
\end{itemize}
Before either netlist is sent to the browser, \code{netlist_for_view} reduces it to what a drawing needs: the top module's ports (direction and bit numbers) and its cells (type, pin directions and connections). Attributes, source locations and net names are dropped, because they carry the server's temporary folder paths and most of the size, and cells get short neutral ids (\code{c0}, \code{c1}, \ldots). The result is still valid Yosys JSON.

Generic synthesis also reports Yosys' generic (unmapped) cell count, the most common cell types, up to five Yosys warnings (with folder paths reduced to file names), flip-flops and latches (with a warning when a latch was inferred, usually an incomplete \code{if} or \code{case}), and a relative area proxy (cells $\times$ 5), with explicit flags \code{liberty_mapped: false}, \code{area_units: "relative_proxy"} and \code{tech_node: null}. When the timing stage produced a real mapped area, the Results tab hides the proxy and shows the mapped area in \umsq{} instead (Chapter~\ref{ch:timing}).

\section{UVM export (\texttt{uvm\_templates.py})}

\code{generate_uvm_bundle(spec)} writes 16 files from templates: interface, sequence item, sequencer, sequence, driver, monitor, agent, scoreboard, coverage, environment, test, package, top module, Makefile, file list and README. It detects clocks, resets and opcode ports to adapt the templates. The top module uses the design's own clock and reset names and the reset's real polarity (\code{rst_n} and \code{reset_b} are active-low, \code{rst} active-high), and the monitor waits for the reset to be \emph{released}. A clock input makes the design sequential even if the specification never uses the word. Only a design that says it is an ALU (or names the operations) and has an operation port of at least 3 bits gets the ALU reference model; for other designs the scoreboard is a placeholder that logs transactions and leaves a marked hook for the user's own reference model. The bundle is zipped and downloaded by the browser. It is export-only and has never been compiled in Questa here.

\begin{warnbox}[What the last review found in the export]
Every generated sequence declared helper tasks \emph{inside} the \code{body()} task, which SystemVerilog does not allow, so no bundle could compile; mux-like specifications crashed the generator; an active-high reset was wired to a port that did not exist; and a 2:1 mux was treated as an ALU (with an illegal enumeration). All four are fixed and tested structurally. The counter scoreboard's reference model is still simplistic, and the bundle is still never simulated here.
\end{warnbox}

\section{Persistence (\texttt{db.py}) and cleanup}

One SQLite table, \code{projects}, with columns for id, prompt, module name, spec, RTL, testbench, explanation, status, pass/fail counts, iterations, creation time, a waveform-truncated flag, and \code{data} (the full response as JSON). \code{save_project}, \code{get_project}, \code{list_projects}, \code{count_projects} and \code{delete_project} do what their names say; the project list also returns each run's model and whether it was a built-in demo run, read from the saved JSON. \code{cleanup_old_workspaces} deletes, at start-up, only folders whose names look like the app's own (letters, digits, \code{_} and \code{-}) and whose \emph{newest file} is older than 24 hours (a folder's own date does not change when a file inside is rewritten, so an active folder once looked old). The workspace is one absolute folder next to the backend code; tests point it at a temporary folder before anything is imported, so a test run never touches your saved projects.

\section{Data contracts (\texttt{models.py})}

\begin{center}
\small
\begin{tabularx}{\textwidth}{@{}lX@{}}
\toprule
\textbf{Model} & \textbf{Holds} \\
\midrule
\code{PortSpec} & name, direction (input/output/inout), width 1--1024 (accepts strings like ``4''; a parameter name defaults to 1), description. \\
\code{RTLDesignSpec} & module name, parameters, ports, behaviour, constraints. \\
\code{RunRequest} & prompt ($\le$2000 chars), target MHz (1--10000), self\_correct, max\_iterations (0--10, default 5), timeout (1--120 s, default 30), model, testbench model (both pattern-checked, Section~\ref{sec:modelid}), an optional confirmed specification (bounded by \code{check_user_spec}), reasoning effort. \\
\code{SpecParseRequest} & prompt, model and reasoning effort for \code{/api/spec/parse}. \\
\code{SimulationRequest} & design id (pattern-checked), RTL and testbench ($\le$50000 chars each), timeout. \\
\code{SimulationResult} & status, counts, coverage, errors, waveform/transcript paths, log excerpt, check rows and groups, \code{measured_toggle}. \\
\code{ToggleCoverage} & percent, bits toggled and total, signals, up to 12 untoggled signals, a partial flag. \\
\code{StrengthReport} & status (ok, partial, skipped), mutants run, caught, invalid, survivors (description and line), note. \\
\code{TimingReport} & available, library, clock, target and period, slack, met, fmax, critical path (start, end, delay and the pin-by-pin rows), area, sequential area, cells, flip-flops, cell types, \code{equivalence} (proven, different, not\_proven or empty) with \code{equivalence_method}, note, error. \\
\code{SecondOpinion} & status (PASS, FAIL, unavailable), counts, the model, a note. \\
\code{IterationRecord} & one pass of the loop (see Chapter~\ref{ch:pipeline}). \\
\code{Waveform}, \code{Schematic} & data for the two visual tabs. \\
\code{RunResponse} & everything above for one run, plus synthesis, \code{fallback_notice}, \code{verification_warning}, \code{model}, \code{testbench_model}, \code{independent_testbench}. \\
\code{UVMExportRequest/Response} & prompt and optional module name; the files, count and base64 zip. \\
\bottomrule
\end{tabularx}
\end{center}

% ---------------------------------------------------------------------------
\chapter{The Backend API}
\label{ch:api}

\begin{center}
\small
\begin{longtable}{@{}>{\ttfamily}l>{\ttfamily}p{4.1cm}p{8.1cm}@{}}
\toprule
\textrm{\textbf{Method}} & \textrm{\textbf{Path}} & \textbf{Purpose} \\ \midrule \endhead
GET & / & Name, provider, offline flag, simulator availability. \\
GET & /health & \texttt{\{"status":"ok","simulator":true\}} for health checks. \\
GET & /api/models & The Model menu: models from all sources, which sources are available (with hints), how many are hidden, and whether the app is offline. \code{?include_paid=true} also lists models not marked free; \code{?refresh=true} re-reads the catalogues. \\
POST & /api/spec/parse & Turn a prompt into a specification and stop, for review before generation. \\
POST & /api/design/stream & The main route: the full pipeline as Server-Sent Events. Saves the finished project. \\
POST & /api/design/run & The same pipeline, returning one JSON response at the end. Saves the project. \\
POST & /api/simulation/run & Simulate a hand-edited RTL/testbench pair once (the Re-run button). No auto-fix. \\
POST & /api/agent/ask & The debugging agent: re-simulates the editor's code, then streams the evidence steps and one answer as Server-Sent Events (Chapter~\ref{ch:agent}). \\
POST & /api/synthesis/run & Generic Yosys synthesis; returns \texttt{\{"available": false\}} when not installed. \\
POST & /api/timing/run & The timing and area estimate on its own (RTL, module name, target, optional testbench for parameter overrides). \\
GET & /api/projects & List saved projects (limit 1--100, offset). \\
GET & /api/projects/\{id\} & Load one saved project. \\
DELETE & /api/projects/\{id\} & Delete one saved project. \\
POST & /api/uvm/export & Generate the UVM bundle (files and zip). \\
POST & /api/admin/cleanup & Delete old workspaces; disabled unless \code{ADMIN_TOKEN} is set and sent as \code{X-Admin-Token}. \\
\bottomrule
\end{longtable}
\end{center}

\textbf{Design choices in \file{main.py}:}
\begin{itemize}
  \item \textbf{CORS: this machine only.} Browsers may call the API only from local addresses (\code{http://localhost:5173}, \code{http://127.0.0.1:5173}, the preview ports 4173 and \code{http://localhost:3000}). Extra origins can be added one by one, as exact addresses, with \code{CORS_ALLOW_ORIGINS}; a wildcard is ignored. There is deliberately no pattern such as \code{*.vercel.app}: anyone can host a page there, and that page could make a visitor's browser call a backend that holds model keys. Credentials (cookies) are not allowed, and only GET, POST, DELETE and OPTIONS are accepted.
  \item Every route except \code{/}, \code{/health} and a plain \code{/api/models} calls \code{_check_rate_limit}: 20 requests per 60 seconds per IP, protected by a lock because handlers run in threads. By default the IP is the real socket address. Only when \code{TRUST_PROXY_HEADERS=1} (behind a trusted proxy such as Render) is the \emph{last} \code{X-Forwarded-For} entry used, which the proxy adds; otherwise any caller could fake that header to dodge the limit.
  \item Simulation, synthesis, timing, model list, specification, project and export routes are plain \code{def}, so FastAPI runs them in a thread pool and the server stays responsive during a simulation or a slow model listing.
  \item Model ids in every request are checked against the safe pattern (Section~\ref{sec:modelid}) before anything else happens; a bad one is HTTP 422.
  \item Errors are logged in full on the server, but clients only receive a short, safe message such as ``Internal server error'', so stack traces and file paths never leak. A storage error says ``The server could not read or write its working files'' with the system's reason, without the path; absolute server paths are removed from short messages, while model ids that look like paths (\code{nim/poolside/...}) are kept.
  \item Request bodies are limited to 1 MB by a small middleware (\code{BodySizeLimit}) that counts the bytes as they arrive, so a chunked upload without a size header is capped too; a larger body gets HTTP 413 before it is parsed.
  \item The server listens on 127.0.0.1 only: it has no login, so it must not be reachable from the room's Wi-Fi. The admin token is compared as bytes, so a non-ASCII header is a clean 403, not a 500.
  \item Design ids are full UUIDs; an earlier 8-character id could collide and overwrite another project.
  \item If a built frontend exists in \file{static/} (the Docker image), the backend serves it too, so one container runs the whole app.
\end{itemize}

% ---------------------------------------------------------------------------
\chapter{The Frontend}
\label{ch:frontend}

\begin{figure}[H]
\centering
\resizebox{\textwidth}{!}{%
\begin{tikzpicture}[node distance=6mm and 5mm,font=\small]
  \node[hbox] (app) {App.tsx \\ \scriptsize all state + handlers};
  \node[box,below left=10mm and 14mm of app,align=center] (pp) {PromptPanel \\ \scriptsize ModelPicker, TestbenchModelSelect, AgentChat, DebugChat};
  \node[box,below=22mm of app] (spec) {SpecReview};
  \node[box,right=24mm of app] (rp) {ResultsPanel};
  \node[box,above right=4mm and 6mm of rp,align=left,font=\scriptsize] (tabs) {CodeEditor (Design) \\ TestCasesView \\ WaveformViewer \\ SchematicView \\ SpecSummary};
  \node[box,left=22mm of app] (api) {api.ts \\ \scriptsize HTTP + SSE};
  \foreach \t in {pp,spec,rp} \draw[arr] (app)--(\t);
  \draw[arr] (rp)--(tabs); \draw[arr,<->] (app)--(api);
\end{tikzpicture}}
\caption{Component tree. App owns the data; components receive it as props and report user actions back.}
\end{figure}

\section{App.tsx: state and handlers}

\code{App} keeps all state with React's \code{useState}: the provider badge, the model list and sources, the chosen model, testbench model and thinking level, the review-the-specification switch, the RTL and testbench text, module name, explanation, result, waveform, schematic, iteration history, design id, synthesis, strength, timing, second opinion, the models that wrote the design and testbench, the specification used, lint info, fallback notice and warning, timeline messages and saved projects. On first load it calls \code{getModels} (which also tells it whether the backend is live or offline) and \code{listProjects}. If the backend cannot be reached, the header shows a red ``Server not reachable'' badge with a Retry button: an unreachable server is not a mode, and the old offline-demo badge made it look like one. The chosen model, testbench model, ``Show other models'' and the review switch are remembered in the browser (\code{localStorage}).

\begin{description}
  \item[\code{loadModels}, \code{chooseModel}] Read the Model menu from \code{/api/models}; keep the current choice if it is still offered, otherwise restore the remembered one, otherwise use the server's default. A testbench model that is no longer offered is cleared.
  \item[\code{handleGenerate}] If the review switch is on (and the app is live), call \code{parseSpec} and show the SpecReview panel; otherwise call \code{startRun} straight away.
  \item[\code{startRun}] Clears \emph{every} result of the previous run, including the RTL, testbench, explanation and design id (otherwise a new run that failed after its RTL arrived left the new RTL next to the old testbench, and Re-run checked that mixed pair), closes any open specification review, opens the stream (with the confirmed specification if there is one), and for each event: adds a timeline message, shows a fallback notice once (\code{lastNoticeRef} prevents duplicates), clears old lint when a repair changes the code, and fills whichever fields the event carries. On \code{done} it loads the complete response (including \code{pass_disputed}), switches to the Results tab and refreshes Projects. It turns network problems into plain messages that say the run is lost and what to do.
  \item[\code{handleCancel}] Aborts the stream or a re-run (the Stop button). Both now carry an abort signal; Stop used to do nothing during a Re-run.
  \item[\code{handleReRun}] Sends the edited code to \code{/api/simulation/run} once, without the AI. It clears strength, synthesis, timing, second opinion, lint and ``Design written by'', because those described the previous code, and shows the note ``Re-run of edited code: only the simulation was repeated \ldots''.
  \item[\code{handleEvidenceRun}, \code{handleApplyFix}] For the debugging agent: when its evidence run simulated code that differs from the results on screen, show that run instead (with a timeline message); put a fix from the agent into the editor and open the Design tab.
  \item[\code{handleUVMExport}] Requests the bundle and downloads the zip; a second click while an export is running is ignored.
  \item[\code{handleLoadProject}, \code{handleDeleteProject}] Open or delete a saved run (deleting asks for confirmation). Loading resets the repair count and lint too, and a counter (\code{loadSeq}) makes sure only the \emph{latest} load writes the screen: two quick clicks used to let the slower, older project overwrite the newer one.
\end{description}

\section{api.ts and streaming}
\label{sec:apits}

The browser's built-in \code{EventSource} cannot send a POST body, so \code{streamDesign} uses \code{fetch} and reads the response body as a stream. It decodes chunks into a buffer, normalises \texttt{\textbackslash r\textbackslash n} to \texttt{\textbackslash n} on the whole buffer (a pair can be split across two chunks), splits on blank lines, takes the \code{data:} line of each frame and parses its JSON. After Stop, frames already received are not applied (a cancelled run used to apply its final result), and the stream is always closed, even if handling an event throws. A FastAPI validation error (a list of messages) is shown as its messages, not as raw JSON. HTTP 502/504 errors say that the server did not answer in time. \file{api.stream.test.ts} runs this real function against a fake server; every other App test replaces it with a stub.

\section{The components}

\begin{description}
  \item[PromptPanel] The sidebar: the ``Describe the block'' text box with a character count, the Generate RTL button, the run timeline (passed in as a child), four example chips (including ``Buggy counter'', scripted only in offline mode), the Model menu, ``Testbench written by'', ``Review the specification before generating'', a thinking-level selector for models that have one, and the Target clock slider (10--500 MHz). Ctrl/Cmd+Enter submits.
  \item[ModelPicker] Three source tabs (OpenCode, Cloudflare, NVIDIA NIM) with arrow-key navigation, a refresh button, the model drop-down grouped by provider (free models marked ``(free)''), the model's note, a hint when a source is not set up, and on the OpenCode tab the ``Show other models'' check box with the number hidden.
  \item[TestbenchModelSelect] ``Testbench written by'': ``Same model as the design'' or any model from any source.
  \item[SpecReview] The specification review panel: module, ports and parameters, editable behaviour and constraints, ``Generate with this specification'' and Cancel.
  \item[AgentChat] The run timeline: one message per event with an icon for its stage and status, the structured diagnosis folded under a ``Diagnosis'' line (it is also on the Results tab), with \texttt{`code`} from the model shown as code, a ``Working\ldots'' timer while a run is going, and the design explanation.
  \item[DebugChat] The ``Ask about this run'' panel under the run timeline: questions, the evidence steps, the answer with clickable times, and the fix card with the server's verdict (Chapter~\ref{ch:agent}).
  \item[CodeEditor] Two tabs (RTL and testbench) using CodeMirror with Verilog highlighting, editable, with a copy button. Edited code can be re-simulated with Re-run. It is shown on the Design tab (before any run, the Design tab shows a welcome page with recent projects instead).
  \item[ResultsPanel] Five tabs: Design, Results, Test cases, Waveform, Schematic. \emph{Results}, top to bottom: the built-in-demo tag when the result is not from an AI model, the verdict (status badge, which reads ``PASS, disputed'' in warning colour when the evidence contradicts the PASS; checks, passed, failed, pass rate), the ``Do not trust this PASS yet'' line, warnings, the second opinion, ``Design written by'', then sections for testbench strength, synthesis, timing and area (with the equivalence line, \code{EquivalenceLine}, and the fold-out critical path table, \code{CriticalPath}), toggle coverage, test types, lint, the self-correction loop, errors and the specification used (SpecSummary).
  \item[TestCasesView] Every executed check, with filters, search and a bit-by-bit comparison (Section~\ref{sec:testcases}).
  \item[WaveformViewer] Draws the waveform as SVG: a time ruler, single-bit signals as high/low lines, buses as labelled hex segments, zoom (0.5$\times$ to 4$\times$) and reset, and a notice when data was truncated. No external tool such as GTKWave is needed. Values are formatted by \code{formatValue}, which never fails: Icarus prints a partly unknown nibble in capitals (\code{X0}), and that once threw an error in decimal mode that blanked the whole page. Only a pinned cursor is announced to screen readers, and a cursor from one run is not carried to the next.
  \item[ErrorBoundary] Wraps the results panel: if any part of it fails to draw, that panel shows the error and a ``Try again'' button, and the rest of the app keeps working.
  \item[SchematicView] Three views chosen with buttons: \emph{Block diagram} (the module as a box with input ports on the left, outputs on the right and widths shown as \code{[n:0]}, drawn by the component itself), \emph{RTL cells} and \emph{Gate level} (the two netlists of Section~\ref{sec:schematicnetlists}). A view without a netlist is greyed out, with the reason as a tooltip. The netlist views are laid out by \textbf{netlistsvg} (open source, MIT licence) through \file{netlistRender.ts}, which is \emph{lazy-loaded}: the layout engine is fetched with a dynamic \code{import} only the first time such a view is opened, so the first page load stays fast. The SVG it returns is shown through an \code{<img>} tag from a blob URL, not inserted into the page: an SVG shown as an image cannot run script, whatever names the design's signals carry. Above 400 cells (\code{MAX_DRAWN_CELLS}) nothing is drawn and the tab says why. Zoom (0.5$\times$ to 6$\times$), \emph{Fit to panel} and \emph{Download SVG} work on the drawn view.
\end{description}

% ---------------------------------------------------------------------------
\chapter{The Debugging Agent}
\label{ch:agent}

Everything so far is a \emph{fixed pipeline}: the orchestrator decides which step runs next, and the model only fills in each step. This chapter describes the one part of SiliconScribe where the model decides what to look at. In the app it is the panel \textbf{Ask about this run} in the left sidebar, under the run timeline. You type a question about the code in the editor (``Why did the first check fail?''), and the answer is built from the simulation itself: values read from the waveform at exact times, lines found in the source, and fixes that the server has simulated.

It was added on 24 November 2026, \emph{after} the project report was submitted (deadline 6 November 2026). In the report it appears only as an extension after submission, never as part of the submitted work. The code is on the branch \code{agent-debug}, commit \code{3d05b94}.

\begin{center}
\small
\begin{tabularx}{\textwidth}{@{}>{\raggedright\arraybackslash}p{5.0cm}X@{}}
\toprule
\textbf{File} & \textbf{What it does} \\
\midrule
\file{backend/agent.py} & The tool loop: the system prompt, the eight tools (\code{Toolbox}), the overview, the grounding rule, fix verification and the rule-based fallback. \\
\file{backend/wave_query.py} & \code{WaveDB}: a queryable copy of one run's VCD (exact values at exact times, for any signal). \\
\file{backend/main.py} & \code{agent_ask}, the route \code{POST /api/agent/ask}: the evidence run and the event stream. \\
\file{backend/llm_service.py} & \code{agent_chat} (one model turn, plain text in and out) and \code{uses_offline}. \\
\file{frontend/src/components/}\newline\file{DebugChat.tsx} & The panel: questions, evidence steps, the answer, clickable times and the fix card. \\
\file{frontend/src/api.ts} & \code{askAgent}, which reads the stream with the same \code{readSSE} function as a normal run. \\
\bottomrule
\end{tabularx}
\end{center}

\section{What an agent is}

A plain chat model answers from the text it was given. If you paste a Verilog file and ask ``why does the hold check fail?'', it reads the code and \emph{guesses} what the hardware did. It cannot see the waveform, so it can describe a waveform that never happened (Section~\ref{sec:llmalonefails} explains why that is the core problem of using a model for hardware).

An \textbf{agent} is a model placed in a loop with \textbf{tools}. A tool is a small function on the server that does one real thing, such as ``give me the value of \code{count} at 25 ns''. In each turn the model either asks for one tool or gives its final answer. The server runs the tool, adds the result to the conversation, and asks the model again. The model chooses the questions; the server supplies the answers from real data.

\begin{analogy}
A doctor on the phone can only guess from what you say. A doctor in a clinic can order a blood test, read the result, then order the next test that result suggests. The tools are the lab tests: the doctor chooses them, but the numbers come from the lab, not from the doctor's memory.
\end{analogy}

\begin{figure}[H]
\centering
\resizebox{\textwidth}{!}{%
\begin{tikzpicture}[node distance=6mm and 9mm]
  \node[hbox] (q) {Your question\\ \scriptsize + the code in the editor};
  \node[box,right=of q] (sim) {Evidence run\\ \scriptsize Icarus, sandbox, 30 s};
  \node[box,right=of sim] (ov) {Overview\\ \scriptsize verdict, failing checks,\\ \scriptsize clocks, signals, RTL};
  \node[box,below=13mm of ov] (m) {Model\\ \scriptsize one JSON object per reply};
  \node[box,left=30mm of m] (tb) {Toolbox\\ \scriptsize 8 tools: waveform, code, test\_fix};
  \node[okbox,below=13mm of m] (ans) {Answer\\ \scriptsize text, suspect, confidence, fix};
  \node[box,left=of ans] (fx) {Server re-simulates the fix\\ \scriptsize its own verdict is shown};
  \node[hbox,left=of fx] (ui) {Panel in\\ the browser};
  \draw[arr] (q)--(sim); \draw[arr] (sim)--(ov); \draw[arr] (ov)--(m);
  \draw[arr] ([yshift=1.5mm]m.west) -- node[above,font=\scriptsize]{tool call} ([yshift=1.5mm]tb.east);
  \draw[arr] ([yshift=-1.5mm]tb.east) -- node[below,font=\scriptsize]{result added} ([yshift=-1.5mm]m.west);
  \draw[arr] (m)--node[right,font=\scriptsize]{answer}(ans);
  \draw[arr] (ans)--(fx); \draw[arr] (fx)--(ui);
  \node[note,right=4mm of m,text width=3.2cm] {at most 10 model calls;\\ no waveform query yet:\\ sent back once};
\end{tikzpicture}}
\caption{One question. The model chooses the tools; every tool result comes from the simulator or the source, and every fix is simulated again by the server before it is shown.}
\label{fig:agentloop}
\end{figure}

\section{The protocol: one JSON object per reply}

The system prompt (\code{SYSTEM_PROMPT} in \file{agent.py}) tells the model that every reply must be exactly one JSON object, in one of two shapes:

\begin{lstlisting}[style=plain]
{"thought": "one short sentence", "tool": "cycle_table",
 "args": {"signals": ["enable", "count"], "from_ns": 0, "to_ns": 100}}

{"answer": "<markdown>", "suspect": "design" | "testbench" | "both" | "unclear" | "none",
 "confidence": "high" | "medium" | "low",
 "fix": {"file": "rtl" | "testbench", "code": "<the COMPLETE corrected file>"}}
\end{lstlisting}

The first is a tool call, the second the final answer (\code{fix} is optional). The prompt also gives the working rules: start from the first failing check, remember that the testbench can be wrong too, test a fix before claiming it works, query the waveform at least once, and write every time as ``@125ns'' so the user can click it.

\subsection{Why plain JSON text and not native function calling}

Some AI services offer \textbf{native function calling}: the program declares its tools in a special field of the request, and the service returns a structured tool call instead of text. SiliconScribe does not use it. Its models come from four kinds of source (the OpenCode app, Cloudflare Workers AI, NVIDIA NIM and any OpenAI-compatible endpoint), and they do not all offer that field in the same way. The OpenCode app, for example, is driven through its command line (Section~\ref{sec:routing}): the prompt goes in as text and the answer comes back as text, with no place to declare tools. Plain JSON in plain text works the same through every source, so the agent works with any model in the Model menu.

The price is that the reply must be parsed. \code{_parse_reply} uses the same forgiving reader as the rest of the app (\code{_extract_json}, Section~\ref{sec:llmparse}), which takes the last complete JSON object in the reply. A reply that is not one of the two shapes counts as a \textbf{format slip}: the first time, the model is told ``Your reply was not one JSON object'' and asked again; after a second slip its text is shown, with the note ``The model did not follow the tool protocol; this text is unchecked.''

Each model turn is one call of \code{agent_chat} in \file{llm_service.py}: the chosen model, temperature 0.1, the system prompt and a transcript that holds the question, up to 4 earlier questions and answers of the conversation, the overview, and every step so far. The last 5 tool results are kept whole (up to 3000 characters each); older ones are shortened to 400 characters, because each turn sends the whole growing transcript again.

\section{The evidence run}

The question is about the code \emph{in the editor}, which may have been edited since the last run. So the first thing \code{agent_ask} in \file{main.py} does is simulate that code again, with \code{_rerun}, the same function the Re-run button uses: a fresh workspace folder (\code{new_workspace_id}), Icarus Verilog, the full sandbox policy and allow-list (Section~\ref{sec:sandbox}), the resource caps, and a 30-second limit. \code{_simulate_checked} also applies the rule that a testbench without checks can never PASS. A \code{$system} call in the code is refused here exactly as in a normal run.

The result is streamed to the browser as a \code{run} event. If the editor's code differs from the code whose results are on screen, the Results and Waveform tabs switch to this run, and the timeline says so (``The code was edited since the last run, so the question re-simulated it \ldots''). That way the times in the answer always match the waveform on screen. If the code is unchanged, nothing on screen changes: Icarus gives the same result for the same code.

\code{overview} then writes a short summary that opens the model's first turn: the verdict and counts; for ERROR or TIMEOUT the error lines and the end of the log; the first six failing checks with their times; how long the simulation ran and the length of one VCD tick; up to two clocks with their period and first rising edge; the list of signals (up to 80, the design's own first); the RTL with line numbers (up to 400 lines); and the length of the testbench (the model reads it with \code{read_code} when it needs to).

\section{WaveDB: reading a VCD for questions}
\label{sec:wavedb}

A VCD file (Section~\ref{sec:vcdbasics}) is a list of value changes in time order. The waveform viewer gets a compact copy (\code{parse_vcd}: at most 40 signals and 2000 changes each, in hex). The agent needs something else: the exact value of \emph{any} signal at \emph{any} time, without ever putting the whole file into a prompt. \code{wave_query.load} builds that copy, called a \code{WaveDB}:

\begin{enumerate}
  \item Read at most 8 MB of the file (the same cap as the viewer; if the file was larger, \code{partial} is set and the overview says late changes are missing), the timescale, and every \code{$var} declaration. Parameters, events and strings are skipped.
  \item Walk the changes. \texttt{\#t} sets the current time; \code{1!} or \code{b0101 code} is a change of one signal. Each value is stored as binary text padded to the signal's width (\code{x} and \code{z} bits are kept). A repeated dump of the same value is not a change.
  \item If a signal changes several times in one time step, only the \textbf{settled} value is kept. That is what the step ends with, and what any later code sees.
  \item Each signal ends up with two lists: the times of its changes and the values. Up to 2000 signals and 200,000 changes per signal are kept.
\end{enumerate}

Looking up a value is then a \textbf{binary search} in the list of times (Python's \code{bisect}), the way you find a word in a dictionary by opening it in the middle rather than reading from page one. \code{resolve} finds a signal by full path (\code{testbench.dut.count}), by the end of a path (\code{dut.count}) or by bare name (\code{count}); when several match, the design's own copy wins over a testbench variable of the same name.

\subsection{value\_at and value\_before: why ``just before the edge''}

Every signal has two lookups, and the difference between them is the most important idea in this chapter:
\begin{description}
  \item[\code{value_at(t)}] the value \emph{after} every change at time $t$: the last change at or before $t$.
  \item[\code{value_before(t)}] the value \emph{just before} $t$: the last change strictly before $t$.
\end{description}

Why both? Take a counter at a rising clock edge at 25 ns. In the VCD, \code{count} has a change \emph{at} 25 ns, because the flip-flops updated on that edge. But the flip-flops decided what to store from the values their inputs had \emph{just before} the edge: at the instant the clock rises, a flip-flop takes a snapshot of its input, and only then does its output change (Section~\ref{sec:seqlogic}). If the agent read \code{enable} ``at'' 25 ns, it would see the value after the edge, which may already have been changed by the testbench in the same time step. That is not what the flip-flop saw. So \code{value_before} answers ``what did the flip-flop sample?'' and \code{value_at} answers ``what did it store?''.

\begin{keyidea}
At a clock edge, read the inputs \emph{just before} the edge and the outputs \emph{after} it. A flip-flop samples before and updates after. Reading both at the same instant is the classic off-by-one-cycle mistake in waveform debugging, and it is exactly the mistake the \code{cycle_table} tool is built to avoid.
\end{keyidea}

This was checked by sabotage: making ``just before'' return the same as ``after'' (together with making the fix check always report PASS) made 4 backend tests fail.

\subsection{Times: nanoseconds and ticks}

The VCD stores times as whole \textbf{ticks} of its timescale. The model always speaks in nanoseconds. \code{to_ticks} converts with $\text{ticks} = \text{ns} \times 1000 / \text{tick length in ps}$, and \code{to_ns} converts back. In the runs shown in this guide one tick is 1 ps, so 30 ns is 30,000 ticks. \code{edges} finds rising or falling edges of a 1-bit signal; as in Verilog, a change from \code{x} or \code{z} to 1 also counts as a rising edge. \code{clocks} finds 1-bit signals named like \code{clk} or \code{clock}.

\section{The eight tools}

\code{Toolbox} in \file{agent.py} holds the run being discussed (its \code{WaveDB}, its result, the RTL and testbench) and offers eight tools. Times in the arguments are nanoseconds.

\begin{center}
\small
\begin{tabularx}{\textwidth}{@{}>{\ttfamily}lX@{}}
\toprule
\textrm{\textbf{Tool}} & \textbf{What it returns} \\
\midrule
failing\_checks & The testbench's failed comparisons: number, time, type, name, expected and actual value (up to 25 per call, with an offset to page through more). \\
values\_at & Each signal's value at a time \emph{and} just before it (``\code{count = 4'h2 (2) (was 4'h1 (1))}''). Without a signal list it reads the design's ports. Up to 12 signals per call. \\
changes & Every change of one signal in a time window, and the value entering the window (at most 40 changes; it says when there are more). \\
cycle\_table & One row per clock edge: each signal's value sampled just \emph{before} the edge (what the flip-flops see) and \emph{after} it. Up to 32 rows. The best tool for sequential logic. The clock is found by name when there is one. \\
find\_value & When a signal held a given value, written in decimal (\code{9}), hex (\code{4'h9}, \code{0x9}), binary (\code{4'b1001}), \code{x} or \code{z} (up to 5 intervals). \\
drivers & Source lines that assign a signal (left side of \code{=} or \code{<=}, not a \code{==} comparison), with the enclosing \code{always} or \code{assign} block, plus instance ports connected to it. This is a \emph{text search} with comments and strings blanked, not a full elaboration of the design, and it says so. \\
read\_code & Numbered lines of the RTL or the testbench (at most 120 per call). \\
test\_fix & Simulates a candidate fix (the complete corrected file) with the other file unchanged, and returns the verdict and the first failing checks. At most 3 per question. \\
\bottomrule
\end{tabularx}
\end{center}

Values are written as \code{4'h9 (9)}: hex always, decimal too when no bit is \code{x} or \code{z} (\code{fmt_value}). A bad call (an unknown signal, a missing time, a fourth \code{test_fix}) raises \code{ToolError}; its message goes back to the model as the tool result, with hints such as similar signal names, so the model can correct itself. A crash inside a tool is logged and reported as ``Tool error'' without ending the conversation. A tool result is cut at 3000 characters. For \code{test_fix}, the code argument is replaced by ``\textless N characters\textgreater'' in the events sent to the browser, so a long file is not echoed back.

\section{The grounding rule}

An answer about a run must rest on the run. The four waveform tools and \code{test_fix} are the tools that read the simulation (\code{WAVE_TOOLS}). If the model answers before any of them has given a real result (a call that ended in a tool error does not count), the answer is not shown. Instead the model is told once: ``Before answering, confirm your explanation with at least one waveform query \ldots, then answer.'' If it answers again without one, that second answer is accepted, so a stubborn model cannot loop forever. A run with no waveform at all (it did not compile) is exempt.

The rule came from a live result. With the model \code{opencode/deepseek-v4.1-flash-free}, the hidden counter bug (A4 in Table~\ref{tab:agentlive}) was diagnosed correctly and the fix passed 57/57, but with \emph{zero} tool calls: the model reasoned from the overview alone. It was right that time, but nothing in the answer came from the waveform.

\section{The answer and the fix}

The final answer carries markdown text, a \textbf{suspect} (design, testbench, both, unclear or none) and a \textbf{confidence} (high, medium or low); anything else is replaced by ``unclear'' or ``medium''. If it includes a fix, \code{_verified_fix} decides what happens to it:
\begin{enumerate}
  \item It must name the RTL or the testbench, contain a module and be at most 50,000 characters; otherwise it is dropped.
  \item A ``fix'' identical to the current file is dropped.
  \item The server simulates it (\code{run_fix}, in the same sandbox, other file unchanged). If the model already tested exactly this code with \code{test_fix}, that result is reused, since the same code gives the same verdict.
  \item The answer carries \emph{that} verdict: status and counts, for example ``PASS, 57 of 57 checks''. If the fix cannot even be simulated, it is shown as ERROR with the reason, not hidden.
\end{enumerate}
The model's own claim that the fix works is never shown as the verdict. A testbench fix also gets a warning in the panel: ``This changes the checker itself. A testbench fix can make a wrong design pass: read it before you apply it.''

\begin{warnbox}[What is checked and what is not]
Only the \emph{fix} is checked, by simulation. The \emph{explanation} is model text and is not verified. In one live run (A6 below) the fix passed 558 of 558 checks, but the explanation was only partly right. A passing fix proves that the new code passes this testbench, not that the model understood why the old code failed.
\end{warnbox}

\section{When there is no AI: rule-based evidence}

\code{rule_based_answer} collects the same kind of evidence with fixed rules: the verdict, the first failing check (name, time, expected and actual), the values of the ports at that time, a cycle table of the edges before it, and the lines that assign the design's outputs. It is used in three cases:
\begin{itemize}
  \item no AI model is connected (\code{uses_offline} says the chosen model would be served by the offline demo provider);
  \item the model call fails: a usage limit, a timeout or another provider error (the reason is named);
  \item the model used all 10 steps without answering (from step 8 on, the transcript asks it to answer now).
\end{itemize}
The text starts with the reason in bold and ``This is evidence collected by fixed rules, not an AI diagnosis'', and the panel labels it \textbf{Rule-based evidence (no AI)}. It is never presented as the AI's answer: there is no silent swap. This follows the same honesty rule as the main pipeline's fallback (Section~\ref{sec:fallback}).

\section{The event stream}

\code{POST /api/agent/ask} answers with Server-Sent Events, like a normal run, but each event has a \code{type} field instead of \code{stage}:

\begin{center}
\small
\begin{tabularx}{\textwidth}{@{}>{\ttfamily}lX@{}}
\toprule
\textrm{\textbf{type}} & \textbf{Meaning and fields} \\
\midrule
status & Progress text: ``Simulating the code in the editor to collect evidence\ldots''. \\
run & The evidence run: \code{result}, the full simulation result with its waveform. \\
tool & One tool call: \code{tool}, \code{args}, the model's \code{thought}, \code{ok} (false for a tool error), a one-line \code{summary} and the \code{result} text. \\
answer & The end: \code{text}, \code{suspect}, \code{confidence}, \code{fix} (with the server's verdict, or null), \code{rule_based}, \code{steps}, the \code{model} (null for rule-based), and \code{note} when the text is unchecked. \\
error & The end, when something failed: a short, safe \code{message}. \\
\bottomrule
\end{tabularx}
\end{center}

Exactly one \code{answer} or \code{error} ends every stream. The request is checked by Pydantic (\code{AgentRequest}): RTL and testbench of 1 to 50,000 characters each, a question of 1 to 2000 characters, at most 6 earlier turns, a model id that matches the safe pattern (Section~\ref{sec:modelid}), and an optional thinking level. The rate limit applies, as on every other route.

\section{The panel in the browser}

\code{DebugChat} is shown in the sidebar once there is an RTL and a testbench. It is disabled while a generation or a Re-run is in progress, and a new design starts a fresh conversation. Its parts:
\begin{description}
  \item[Asking] Three example questions (``Why did the first check fail?'', ``Is the design or the testbench wrong?'', ``Walk me through the first 5 clock cycles.''), or your own question (Enter sends, Shift+Enter adds a line). \code{ask} sends the code exactly as it is at that moment, the chosen model and thinking level, and up to 4 earlier questions and answers. \textbf{Stop} aborts the request.
  \item[Evidence steps] Each \code{tool} event becomes one numbered step with a plain label (``Built a cycle-by-cycle table: enable, count 0--100ns''). Click a step to see the model's thought and the raw tool result. A failed tool call is marked.
  \item[The answer] A header that says \textbf{AI answer} with the number of evidence steps, or \textbf{Rule-based evidence (no AI)}; the model id; ``Suspect: Design \textbullet\ confidence high'' (hidden when unclear); the note, if any; then the text.
  \item[Plain text only] \code{RichText} shows the model's text as plain text. Only three things are recognised: \texttt{`code`}, \texttt{**bold**} and times. Nothing is ever inserted as HTML, so a model cannot put a script or a fake button into the page.
  \item[Clickable times] A time such as ``@30ns'' or ``30 ns'' becomes a button. The waveform viewer works in VCD ticks, so \code{inline} converts: $\text{ticks} = \text{round}(\text{ns} \times 1000 / \text{timescale\_ps})$. Clicking it opens the Waveform tab with cursor A at that time. A sabotage test proved this matters: a wrong ns-to-tick conversion made a frontend test fail.
  \item[The fix card] Green for PASS, red otherwise: ``Proposed fix to the design: the server re-simulated it, PASS (57 of 57 checks passed)'', the testbench warning when the fix is to the testbench, the failing checks when it did not pass, a fold-out with the proposed code, and \textbf{Apply to the editor}. Applying puts the code in the editor and opens the Design tab; the main results change only when you press \textbf{Re-run}.
\end{description}

When the \code{run} event arrives, \code{handleEvidenceRun} in \file{App.tsx} compares the code that was asked about with the code whose results are on screen; if they differ, the evidence run replaces the results, exactly as a Re-run would (\code{applyRerun}). \code{askAgent} reads the stream with the same \code{readSSE} function as \code{streamDesign} (Section~\ref{sec:apits}), treats \code{answer} and \code{error} as final, and says ``The connection closed before the answer arrived'' if neither came.

\section{A worked example: the hidden counter bug}
\label{sec:agentexample}

The built-in 4-bit counter (Section~\ref{sec:counterexample}) was given a bug with no comment pointing at it: \code{else if (enable)} became \code{else if (enable || count[0])}. Whenever the lowest bit is 1, the counter now counts even with \code{enable} low. The testbench is unchanged; its hold checks set \code{enable} to 0 and expect \code{count} to stay put. The simulator says FAIL, 12 of 57 checks passed (45 failed).

On 24 November the question was asked live with the model \code{nvidia/nemotron-3-ultra-550b-a55b} on NVIDIA NIM (A2 in Table~\ref{tab:agentlive}). The answer took 34 s. The model built a cycle table, which showed \code{count} changing while \code{enable} was 0, at 25, 55 and 85 ns. This is what that table looks like (reproduced afterwards by running the same tool on the same code; the rows from 15 to 55 ns):

\begin{lstlisting}[style=plain]
posedge of testbench.clk; 'before' = sampled by flip-flops at the edge, 'after' = value once
the edge's updates settled.
cycle | time | enable before -> after | count before -> after
    1 | 15ns | 1                      | 4'h0 (0) -> 4'h1 (1)
    2 | 25ns | 0                      | 4'h1 (1) -> 4'h2 (2)     <- enable was 0, count moved
    3 | 35ns | 0                      | 4'h2 (2)                 <- count[0] is 0: it holds
    4 | 45ns | 1                      | 4'h2 (2) -> 4'h3 (3)
    5 | 55ns | 0                      | 4'h3 (3) -> 4'h4 (4)     <- moved again
\end{lstlisting}

Read row 2 with the key idea above: just before the edge at 25 ns, \code{enable} was 0 and \code{count} was 1; after it, \code{count} was 2. A correct counter must hold there. Row 3 is the clue to the cause: with \code{count} = 2 (lowest bit 0) it does hold. So the counter moves only when the lowest bit is 1, which points straight at \code{count[0]} in the condition. The testbench's first hold check, at 30 ns, expected 1 and got 2.

The answer correctly blamed the condition in the design, and its proposed fix was re-simulated by the server: \textbf{PASS, 57 of 57}. Figures~\ref{fig:askpanel} and~\ref{fig:askjump} in Chapter~\ref{ch:usage} show the same bug asked in the real browser (A3).

\section{Live results}

These are single runs on small designs, mainly with one model, on 24 November 2026: evidence that the mechanism works, not a benchmark. The model is \code{nvidia/nemotron-3-ultra-550b-a55b} on NVIDIA NIM unless stated.

\begin{center}
\small
\begin{longtable}{@{}l>{\raggedright\arraybackslash}p{3.9cm}>{\raggedright\arraybackslash}p{6.1cm}>{\raggedleft\arraybackslash}p{2.2cm}@{}}
\caption{Live results of the debugging agent.}\label{tab:agentlive}\\
\toprule
\textbf{\#} & \textbf{Case} & \textbf{Result} & \textbf{Time} \\ \midrule \endfirsthead
\toprule
\textbf{\#} & \textbf{Case} & \textbf{Result} & \textbf{Time} \\ \midrule \endhead
A1 & Built-in buggy counter (adds 2; the RTL has a ``BUG'' comment) & Correct, design; fix re-simulated 57/57 & 31 s \\
A2 & Hidden bug \code{enable || count[0]}, no comment & Correct; used a cycle table (count moves while enable = 0 at 25, 55, 85 ns); fix 57/57 & 34 s \\
A3 & Same hidden bug, in the real browser (twice) & Correct both times, fix 57/57; one answer cited 9 clickable times & 107 s, 74 s \\
A4 & Same hidden bug, model \file{opencode/deepseek-v4.1-flash-free} & Correct, fix 57/57, but zero tool calls (reasoned from the overview): led to the grounding rule & 49 s \\
A5 & Correct counter, wrong testbench (expects 3 instead of 2 in \code{counting_continues}) & Suspect \textbf{testbench}: the check runs at the falling edge at 560 ns, before the third rising edge; testbench fix 57/57 & 219 s, 5 steps \\
A6 & Saved SR-latch project that timed out (a real AI design from 20 Nov), real browser & Fix (gate delays) re-simulated 558/558, but the explanation was only partly right: it blamed the hold state; the real cause is the testbench releasing both inputs from 0,0 at once (Section~\ref{sec:latch-timeout}) & 316 s \\
A7 & Older saved counter whose 20-check testbench never drops enable, with the hidden bug & Said the run passed and no hold check failed (the question's premise was false) instead of inventing a failure; this weak testbench does not expose the bug & about 170 s \\
\bottomrule
\end{longtable}
\end{center}

Three of these teach the most. \textbf{A5}: the agent can say ``your design is fine, the checker is wrong'', with evidence (the time of the check relative to the clock edges). \textbf{A6}: the fix is verified, the explanation is not. \textbf{A7}: when the question assumes something false, the evidence lets the model say so.

\begin{warnbox}[Limitations]
\begin{itemize}
  \item The explanation is model text; only fixes are verified by simulation (A6).
  \item It is slow: about 30 s to 5 minutes per question on free models, because each step sends the growing transcript again.
  \item \code{drivers} is a text search, not an elaborated netlist. Icarus only: no SVA, no UVM.
  \item One run per case, small designs, mainly one model: not a benchmark.
  \item The evidence run is limited to 30 s, so a design that hangs (A6) costs the full 30 s before the model even starts.
\end{itemize}
\end{warnbox}

\begin{keyidea}
The agent does not weaken ``only the simulator decides''. The model now chooses \emph{which} questions to ask, but every answer to those questions comes from the simulator's waveform or the source text, and every fix is simulated again by the server. What the model adds is the reasoning that connects the evidence, and that part is labelled as model text.
\end{keyidea}

\begin{viva}
\emph{``Isn't this just ChatGPT with your code pasted in?''} No. A chat model guesses what the waveform looked like. Here the model must ask for the waveform: values at exact times, a table of every clock edge with the inputs sampled just before the edge, the lines that drive a signal. An answer without at least one simulation query is sent back once. And a proposed fix is never shown with the model's claim; the server simulates it and shows its own verdict.
\end{viva}

% ---------------------------------------------------------------------------
\chapter{Lessons from the Nine-Role Review}
\label{ch:review}

Before the viva the whole tool (not the documents) was reviewed once more, this time by nine reviewers, each acting as one kind of engineer. Each role looks at the same code with different questions, so each finds different defects. This chapter explains what each role looks for, the most instructive thing it found, \emph{why} the defect existed, and why the fix works. It is the best preparation for the question ``how do you know it works?''.

\section{How the review was run}

\begin{enumerate}
  \item \textbf{Nine roles, read-only.} ASIC design, design verification, backend systems, security, LLM integration, frontend, interface design, test engineering and release engineering. Each reviewer could read everything and run experiments in a scratch copy, but could not change the project. About 110 findings were reported.
  \item \textbf{Reproduce first.} No finding was fixed on someone's word. Each was reproduced with a real run (a real simulation, a real Yosys or OpenSTA run, a real browser state in a test) before any code changed.
  \item \textbf{A test for every fix.} Each fix comes with a regression test that fails without it.
  \item \textbf{Sabotage the fixes.} Each fix was then undone on purpose, one at a time, to confirm that a test fails. Four fixes survived this at first: their tests passed for the wrong reason and were rewritten.
  \item \textbf{An independent verifier.} A tenth reviewer received only the original findings and the change set (not the author's explanations), tried to break every security and verdict fix, and re-ran all saved runs with the old and new code. It found a sandbox bypass through a macro argument, a regular-expression slowdown introduced by the fixes, and three regressions. Those went through the same steps.
\end{enumerate}

\begin{keyidea}
Fixing is the easy part. The discipline is in the steps around it: reproduce the defect, pin it with a test that fails, prove the test can fail by undoing the fix, and let someone who did not write the fix try to break it. The independent round showed why: three of the author's own fixes had introduced new bugs.
\end{keyidea}

\section{What each role looks for, and what it found}

\subsection*{ASIC design engineer: ``are the numbers physically right?''}
An ASIC engineer distrusts any number until it is traced to the tool output. Here: Yosys writes large areas as \code{1.6E+03}, and the parser read only plain decimals, so a 64-bit register ``mapped to no logic''. Input and output delays were attached to the first clock only, so a slow path on a second clock was reported as meeting 500 MHz. The maximum frequency was computed with a whole-period formula even for half-cycle paths. And the equivalence proof's induction assumes all flip-flops tick together and start equal, so it ``proved'' a netlist with a flip-flop on the wrong clock edge and a design that relied on a power-up value. \emph{Lesson:} every formula has a domain where it is valid; report ``not proven'' outside it.

\subsection*{Design-verification engineer: ``can a wrong design be called PASS?''}
A verification engineer attacks the oracle and the verdict. Found: a later \code{Failed: 0} cancelled an earlier \code{Failed: 4}; \code{CHK|FAIL} rows were ignored; a design could forge the summary from a \code{final} block; the standard check compared only 64 bits; a slow testbench scored 100\% strength because timeouts counted as catches; and a repair could keep the bug and edit the expected values, leaving a green PASS that its own original testbench contradicted. \emph{Lesson:} any evidence against a PASS must win, and must be visible in the headline, not only in a warning paragraph.

\subsection*{Backend systems engineer: ``what happens under concurrency and failure?''}
Found: two runs of the same design shared one folder and swapped verdicts (25 of 25 tries); a compile timeout killed \code{iverilog} but left its child \code{ivl} spinning at 100\% CPU; a runaway testbench made the parser hold every check line in memory (about 470 MB); an old waveform was shown for a new run. \emph{Lesson:} a subprocess is a tree, a folder is shared state, and anything you read must be bounded.

\subsection*{Security engineer: ``can untrusted code reach outside its box?''}
Found: a comment between \code{$dumpfile} and its argument slipped past the path check and Icarus wrote the waveform to any path. Later the verifier found that a macro argument could deliver the \code{$} of \code{$fopen}. Also, the OpenCode child process received every API key in the environment. \emph{Lesson:} check what the tool will \emph{actually} see, which is why the rules now run on the macro-expanded source from Icarus' own preprocessor, and give child processes an allowlist, not a denylist.

\subsection*{LLM integration engineer: ``what does the app do with a bad reply?''}
Found: a Cloudflare model you chose was silently replaced by the demo when the default provider was missing; a quota error at the repair step threw away the whole run; a draft followed by a corrected module produced the module twice; reasoning text was parsed as the answer; error reasons were guessed from substrings (``token'' in \code{max_tokens} read as a credentials problem). \emph{Lesson:} assume replies are messy and services fail mid-run, keep the user's work, and give the real reason.

\subsection*{Frontend engineer: ``can the screen show something false or break?''}
Found: one unusual value (\code{X0}) in decimal mode blanked the whole page; a slow project load overwrote a newer one; Stop did nothing during a Re-run; a new run that failed early left new RTL next to the old testbench, and Re-run checked that mixed pair. \emph{Lesson:} every async result needs an identity, and one broken panel must not take down the app.

\subsection*{Interface designer: ``what will an examiner believe from the screen?''}
Found: the green PASS was louder than the caution that qualified it; saved demo runs reopened looking live; nothing said a live run takes minutes; browser zoom on a laptop projector switched to a stacked layout; the welcome text claimed more than the app proves. \emph{Lesson:} honesty is also visual hierarchy: the most important caveat must sit next to the verdict.

\subsection*{Test engineer: ``would the tests catch a regression?''}
Found by changing one line at a time: the real streaming client had no test (every App test used a stub), the anti-cheat and oscillation tests passed even with those guards removed, the admin endpoint and project deletion had no tests (a change that deleted every project passed the suite), and simply collecting the tests ran the clean-up on the real workspace. \emph{Lesson:} a test you have never seen fail is not yet a test.

\subsection*{Release engineer: ``will the demo run on the day?''}
Found: the backend listened on every network interface without a login; \code{.env.example} had active lines that override the command line; Vite silently moved to another port when 5173 was busy, leaving an old build on the expected address; preflight called a backend left in offline mode ``responding''. \emph{Lesson:} check what is actually running, not that something answers.

\section{After the review: a latch that stopped time}
\label{sec:latch-timeout}

A live run on 20 November, after the review, showed a gap that none of the nine roles had found. The prompt was simply ``design sr latch''. The run ended in \textbf{TIMEOUT} with ``0 checks, 0\%'', which looked like total failure. In fact:
\begin{enumerate}
  \item The model's first RTL used cross-coupled NOR gates, although its own truth table (active-low inputs: hold at 1,1; set at 0,1; reset at 1,0) needs NAND gates. The repair loop diagnosed this correctly and switched to \code{q = ~(s_n & q_n)}, \code{q_n = ~(r_n & q)}, the textbook latch.
  \item The repaired design then passed all 552 checks the testbench printed. After the last one, the testbench drove the forbidden input $s\_n = r\_n = 0$ (both outputs become 1) and released \emph{both} inputs to 1 in the same time step.
  \item In a real chip one gate would be slightly faster and the latch would settle (possibly after a metastable moment). A simulator without gate delays has no ``slightly faster'': both NANDs see 1,1 and switch to 0 together, then both see 0 and switch back to 1, forever. Every switch is a new \emph{delta cycle} at the \emph{same} simulated time, so time never advances.
  \item The testbench's watchdog (\texttt{\#20000; \$finish;}) waits for simulated time, so it never fired. Only SiliconScribe's 30-second wall-clock limit stopped the run.
\end{enumerate}

The verdict was right (the run never finished, so it is not a PASS), but the screen hid the evidence. Three fixes followed, each pinned by a test that fails when the fix is removed:
\begin{itemize}
  \item \textbf{The checks that ran are kept.} On a timeout the simulator now sends \code{SIGTERM} first and gives \code{vvp} two seconds to exit; \code{vvp} flushes its output on \code{SIGTERM}, whereas the old \code{SIGKILL} lost the last few kilobytes (466 of 553 lines survived). The result stays TIMEOUT with no verdict counts, but the header reads ``Checks before stop: 552'', the pass rate is shown as ``---'' (not ``0\%''), and the Test cases tab lists the checks as progress, not a verdict.
  \item \textbf{The likely cause is stated.} The timeout message now says that no watchdog fired, that simulated time therefore stopped or the run needs far longer, and that the usual cause of stopped time is a zero-delay feedback loop (cross-coupled gates, an \code{assign} that feeds itself, an \code{always} block that retriggers itself). The repair model reads the same text, and the testbench prompt now says: leave a forbidden input state one input at a time.
  \item \textbf{A stale explanation is labelled.} The model writes its explanation for the first draft, and repairs do not rewrite it, so the Run panel still described NOR gates. When a repair has changed the code, a note now says the explanation was written for the first draft.
\end{itemize}

\emph{Lesson:} a verdict can be correct and still mislead if the evidence behind it is thrown away. ``0 checks'' and ``552 checks, then the simulator stalled'' are both TIMEOUT, but only the second tells you what to fix.

\begin{viva}
\emph{``Why did the SR latch time out if the design was correct?''} Because the testbench released the forbidden 0,0 input on both pins in the same time step. With zero gate delay the two NAND gates switch together forever in delta cycles, so simulated time stops, and a \texttt{\#delay} watchdog cannot fire; only the wall-clock limit ends the run. Real hardware would settle because one gate is always slightly faster. The fix belongs in the testbench: release one input at a time.
\end{viva}

\section{Numbers to remember}

\begin{center}
\begin{tabularx}{\textwidth}{@{}X>{\raggedleft\arraybackslash}p{5.2cm}@{}}
\toprule
Findings reported by the nine roles & about 110 \\
Each fix pinned by a test that fails without it & yes (sabotage-checked) \\
Defects found by the independent verifier in the fixes & 1 bypass, 1 slowdown, 3 regressions \\
Saved runs with identical verdicts, old vs new code & 123 of 123 (and 14 of 14 evaluation designs) \\
Automated tests after the review & 591 backend + 267 frontend \\
\quad (after the debugging agent, 24 Nov) & 632 backend + 277 frontend \\
\bottomrule
\end{tabularx}
\end{center}

\begin{viva}
\emph{``Did you test your own tool, or only trust it?''} It was reviewed by nine role-based reviewers; every finding was reproduced with a real run, fixed, and pinned by a test that was shown to fail when the fix is removed; then an independent reviewer tried to break the fixes and found three regressions and a bypass, which were fixed the same way. The 123 saved runs give the same verdicts with the old and new code, so the fixes did not change any honest result.
\end{viva}
